
\documentclass{article}
\usepackage{colortbl}
\usepackage{makecell}
\usepackage{multirow}
\usepackage{supertabular}

\begin{document}

\newcounter{utterance}

\centering \large Interaction Transcript for game `clean\_up', experiment `0\_easy\_7obj\_cs', episode 0 with gemma{-}4{-}26b{-}a4b{-}it.
\vspace{24pt}

{ \footnotesize  \setcounter{utterance}{1}
\setlength{\tabcolsep}{0pt}
\begin{supertabular}{c@{$\;$}|p{.15\linewidth}@{}p{.15\linewidth}p{.15\linewidth}p{.15\linewidth}p{.15\linewidth}p{.15\linewidth}}
    \# & $\;$A & \multicolumn{4}{c}{Game Master} & $\;\:$B\\
    \hline

    \theutterance \stepcounter{utterance}  
    & & \multicolumn{4}{p{0.6\linewidth}}{
        \cellcolor[rgb]{0.9,0.9,0.9}{
            \makecell[{{p{\linewidth}}}]{
                \texttt{\tiny{[P1$\langle$GM]}}
                \texttt{Jsem váš herní mistr a hrajete společnou hru s následující mřížkou jako herním plánem:} \\
\\ 
\texttt{```} \\
\texttt{ 1234567} \\
\texttt{╔═╤═══╤═╗} \\
\texttt{║◌│◌DT│◌║ 1} \\
\texttt{╟─┼──┬┘◌║ 2} \\
\texttt{║◌│◌◌│◌◌║ 3} \\
\texttt{╟─┴──┘◌O║ 4} \\
\texttt{║M◌◌P◌◌◌║ 5} \\
\texttt{║◌◌◌◌◌◌◌║ 6} \\
\texttt{║S◌◌◌◌◌A║ 7} \\
\texttt{╚═══════╝} \\
\\ 
\texttt{```} \\
\\ 
\texttt{* Horní okraj zobrazuje souřadnice x rostoucí doprava a pravý okraj souřadnice y rostoucí dolů.} \\
\texttt{* Na vaší mřížce jsou náhodně umístěny následující objekty: 'P', 'O', 'T', 'S', 'D', 'A', 'M'.} \\
\\ 
\\ 
\texttt{Druhý hráč vidí variantu herní desky, kde jsou objekty umístěny na různých náhodných místech. Vy nevidíte desku druhého hráče a on nevidí tu vaši.} \\
\\ 
\\ 
\texttt{**Cíl:**} \\
\\ 
\texttt{Oba hráči musí přesunout objekty na svém pozadí tak, aby stejné objekty skončily na stejných souřadnicích. Musíte komunikovat s druhým hráčem, abyste se dohodli na společné cílové konfiguraci.} \\
\\ 
\\ 
\texttt{**Pravidla:**} \\
\\ 
\texttt{* V každém tahu můžete poslat přesně jeden z následujících dvou příkazů:} \\
\texttt{ 1. `ŘEKNI: <ZPRÁVA>`: odeslání zprávy (vše až do dalšího zalomení řádku) druhému hráči. Přepošlu ji vašemu partnerovi.} \\
\texttt{2. `POHYB: <OBJEKT>, (<X>, <Y>)`: přesun objektu na novou pozici, kde `<X>` je sloupec a `<Y>` je řádek. Budu vás informovat, zda byl váš přesun úspěšný, nebo ne.} \\
\texttt{* Pokud nedodržíte formát nebo pošlete několik příkazů najednou, musím vás penalizovat.} \\
\texttt{* Pokud oba hráči nasbírají více než 16 penalizací, oba prohrajete hru.} \\
\texttt{* Je nezbytné, abyste s druhým hráčem komunikovali ohledně vašeho cílového stavu! *Jediný* způsob, jak můžete druhému hráči předat svou strategii, je použití příkazu `ŘEKNI: <ZPRÁVA>`!} \\
\\ 
\\ 
\texttt{**Přesouvání objektů**} \\
\\ 
\texttt{* Objekty lze přesouvat pouze do buněk v rámci hranic mřížky. Cílová buňka musí být prázdná, tj. musí obsahovat pouze symbol '◌'.} \\
\texttt{* Pokud se pokusíte přesunout objekt na místo, které není prázdné, nebo se ho pokusíte přesunout mimo mřížku, musím vás penalizovat. Dostanete další pokus.} \\
\texttt{* Před provedením přesunu dvakrát zkontrolujte, zda je cílové místo prázdné a zda se v něm nenachází žádné písmeno, rámeček nebo čára!} \\
\\ 
\\ 
\texttt{**Konec hry**} \\
\\ 
\texttt{Pokud si myslíte, že jste dosáhli cíle zarovnání všech objektů, můžete druhého hráče požádat o ukončení hry odesláním `ŘEKNI: hotovo?`. Pokud vás druhý hráč požádá o ukončení hry a vy odpovíte `ŘEKNI: hotovo!`, hra skončí.} \\
\\ 
\texttt{Oba hráči vyhrají, pokud hra skončí během 28 kol, přičemž jedno kolo je definováno jako odeslání platného příkazu oběma hráči.} \\
\\ 
\texttt{**Bodování:**} \\
\\ 
\texttt{Čím blíže jsou stejné objekty na obou herních deskách, tím více bodů získáte. Penalizace vám body snižují. Podaří se vám překonat rekord?} \\
\\ 
\\ 
\texttt{Pošlete prosím zprávu druhému hráči a začněte hru!} \\
            }
        }
    }
    & \\ \\

    \theutterance \stepcounter{utterance}  
    & \multicolumn{4}{p{0.6\linewidth}}{
        \cellcolor[rgb]{0.9,0.9,0.9}{
            \makecell[{{p{\linewidth}}}]{
                \texttt{\tiny{[P1$\rangle$GM]}}
                \texttt{ŘEKNI: Ahoj! Musíme najít společnou cílovou konfiguraci. Já vidím objekty na těchto pozicích: M(1,5), P(4,5), S(1,7), A(7,7), D(5,1), T(6,1), O(7,4). Navrhuji, abychom se pokusili posunout všechny objekty do řádků 2 až 6, aby se to lépe koordinovalo. Jak vypadá tvá deska? Napíš mi pozice svých objektů?} \\
            }
        }
    }
    & & \\ \\

    \theutterance \stepcounter{utterance}  
    & & & \multicolumn{2}{p{0.3\linewidth}}{
        \cellcolor[rgb]{0.9,0.9,0.9}{
            \makecell[{{p{\linewidth}}}]{
                \texttt{\tiny{[GM$|$GM]}}
                \texttt{ŘEKNI: Ahoj! Musíme najít společnou cílovou konfiguraci. Já vidím objekty na těchto pozicích: M(1,5), P(4,5), S(1,7), A(7,7), D(5,1), T(6,1), O(7,4). Navrhuji, abychom se pokusili posunout všechny objekty do řádků 2 až 6, aby se to lépe koordinovalo. Jak vypadá tvá deska? Napíš mi pozice svých objektů?} \\
            }
        }
    }
    & & \\ \\

    \theutterance \stepcounter{utterance}  
    & & \multicolumn{4}{p{0.6\linewidth}}{
        \cellcolor[rgb]{0.9,0.9,0.9}{
            \makecell[{{p{\linewidth}}}]{
                \texttt{\tiny{[GM$\rangle$P2]}}
                \texttt{Jsem váš herní mistr a hrajete společnou hru s následující mřížkou jako herním plánem:} \\
\\ 
\texttt{```} \\
\texttt{ 1234567} \\
\texttt{╔═╤═══╤═╗} \\
\texttt{║◌│◌T◌│◌║ 1} \\
\texttt{╟─┼──┬┘◌║ 2} \\
\texttt{║A│M◌│◌◌║ 3} \\
\texttt{╟─┴──┘◌S║ 4} \\
\texttt{║◌◌◌◌◌◌◌║ 5} \\
\texttt{║◌P◌◌◌◌◌║ 6} \\
\texttt{║◌◌O◌D◌◌║ 7} \\
\texttt{╚═══════╝} \\
\\ 
\texttt{```} \\
\\ 
\texttt{* Horní okraj zobrazuje souřadnice x rostoucí doprava a pravý okraj souřadnice y rostoucí dolů.} \\
\texttt{* Na vaší mřížce jsou náhodně umístěny následující objekty: 'P', 'O', 'T', 'S', 'D', 'A', 'M'.} \\
\\ 
\\ 
\texttt{Druhý hráč vidí variantu herní desky, kde jsou objekty umístěny na různých náhodných místech. Vy nevidíte desku druhého hráče a on nevidí tu vaši.} \\
\\ 
\\ 
\texttt{**Cíl:**} \\
\\ 
\texttt{Oba hráči musí přesunout objekty na svém pozadí tak, aby stejné objekty skončily na stejných souřadnicích. Musíte komunikovat s druhým hráčem, abyste se dohodli na společné cílové konfiguraci.} \\
\\ 
\\ 
\texttt{**Pravidla:**} \\
\\ 
\texttt{* V každém tahu můžete poslat přesně jeden z následujících dvou příkazů:} \\
\texttt{ 1. `ŘEKNI: <ZPRÁVA>`: odeslání zprávy (vše až do dalšího zalomení řádku) druhému hráči. Přepošlu ji vašemu partnerovi.} \\
\texttt{2. `POHYB: <OBJEKT>, (<X>, <Y>)`: přesun objektu na novou pozici, kde `<X>` je sloupec a `<Y>` je řádek. Budu vás informovat, zda byl váš přesun úspěšný, nebo ne.} \\
\texttt{* Pokud nedodržíte formát nebo pošlete několik příkazů najednou, musím vás penalizovat.} \\
\texttt{* Pokud oba hráči nasbírají více než 16 penalizací, oba prohrajete hru.} \\
\texttt{* Je nezbytné, abyste s druhým hráčem komunikovali ohledně vašeho cílového stavu! *Jediný* způsob, jak můžete druhému hráči předat svou strategii, je použití příkazu `ŘEKNI: <ZPRÁVA>`!} \\
\\ 
\\ 
\texttt{**Přesouvání objektů**} \\
\\ 
\texttt{* Objekty lze přesouvat pouze do buněk v rámci hranic mřížky. Cílová buňka musí být prázdná, tj. musí obsahovat pouze symbol '◌'.} \\
\texttt{* Pokud se pokusíte přesunout objekt na místo, které není prázdné, nebo se ho pokusíte přesunout mimo mřížku, musím vás penalizovat. Dostanete další pokus.} \\
\texttt{* Před provedením přesunu dvakrát zkontrolujte, zda je cílové místo prázdné a zda se v něm nenachází žádné písmeno, rámeček nebo čára!} \\
\\ 
\\ 
\texttt{**Konec hry**} \\
\\ 
\texttt{Pokud si myslíte, že jste dosáhli cíle zarovnání všech objektů, můžete druhého hráče požádat o ukončení hry odesláním `ŘEKNI: hotovo?`. Pokud vás druhý hráč požádá o ukončení hry a vy odpovíte `ŘEKNI: hotovo!`, hra skončí.} \\
\\ 
\texttt{Oba hráči vyhrají, pokud hra skončí během 28 kol, přičemž jedno kolo je definováno jako odeslání platného příkazu oběma hráči.} \\
\\ 
\texttt{**Bodování:**} \\
\\ 
\texttt{Čím blíže jsou stejné objekty na obou herních deskách, tím více bodů získáte. Penalizace vám body snižují. Podaří se vám překonat rekord?} \\
\\ 
\texttt{Druhý hráč zahájil hru odesláním této zprávy:} \\
\\ 
\texttt{"Ahoj! Musíme najít společnou cílovou konfiguraci. Já vidím objekty na těchto pozicích: M(1,5), P(4,5), S(1,7), A(7,7), D(5,1), T(6,1), O(7,4). Navrhuji, abychom se pokusili posunout všechny objekty do řádků 2 až 6, aby se to lépe koordinovalo. Jak vypadá tvá deska? Napíš mi pozice svých objektů?"} \\
\\ 
\texttt{Jaký je váš první příkaz?} \\
            }
        }
    }
    & \\ \\

    \theutterance \stepcounter{utterance}  
    & & & \multicolumn{4}{p{0.6\linewidth}}{
        \cellcolor[rgb]{0.9,0.9,0.9}{
            \makecell[{{p{\linewidth}}}]{
                \texttt{\tiny{[GM$\langle$P2]}}
                \texttt{ŘEKNI: Ahoj! Moje pozice jsou: T(4,1), A(1,3), M(3,3), S(7,4), P(2,6), O(3,7), D(5,7). Tvůj návrh na řádky 2{-}6 je dobrý. Navrhuji, abychom se pokusili o jednoduchou konfiguraci, kde každý objekt bude mít svou vlastní souřadnici v řádcích 2 až 6, například: A(1,2), M(2,2), T(3,2), S(4,2), P(5,2), O(6,2), D(7,2). Co na to říkáš?} \\
            }
        }
    }
     \\ \\

    \theutterance \stepcounter{utterance}  
    & & & \multicolumn{2}{p{0.3\linewidth}}{
        \cellcolor[rgb]{0.9,0.9,0.9}{
            \makecell[{{p{\linewidth}}}]{
                \texttt{\tiny{[GM$|$GM]}}
                \texttt{ŘEKNI: Ahoj! Moje pozice jsou: T(4,1), A(1,3), M(3,3), S(7,4), P(2,6), O(3,7), D(5,7). Tvůj návrh na řádky 2{-}6 je dobrý. Navrhuji, abychom se pokusili o jednoduchou konfiguraci, kde každý objekt bude mít svou vlastní souřadnici v řádcích 2 až 6, například: A(1,2), M(2,2), T(3,2), S(4,2), P(5,2), O(6,2), D(7,2). Co na to říkáš?} \\
            }
        }
    }
    & & \\ \\

    \theutterance \stepcounter{utterance}  
    & & \multicolumn{4}{p{0.6\linewidth}}{
        \cellcolor[rgb]{0.9,0.9,0.9}{
            \makecell[{{p{\linewidth}}}]{
                \texttt{\tiny{[P1$\langle$GM]}}
                \texttt{Vaše zpráva byla předána druhému hráči.} \\
\\ 
\\ 
\texttt{Právě hrajete kolo 1 z maximálně 28.} \\
\texttt{Dohromady jste nasbírali 0 z 16 pokut.} \\
\texttt{Druhý hráč poslal tuto zprávu:} \\
\\ 
\\ 
\texttt{"Ahoj! Moje pozice jsou: T(4,1), A(1,3), M(3,3), S(7,4), P(2,6), O(3,7), D(5,7). Tvůj návrh na řádky 2{-}6 je dobrý. Navrhuji, abychom se pokusili o jednoduchou konfiguraci, kde každý objekt bude mít svou vlastní souřadnici v řádcích 2 až 6, například: A(1,2), M(2,2), T(3,2), S(4,2), P(5,2), O(6,2), D(7,2). Co na to říkáš?"} \\
\\ 
\\ 
\texttt{Jaký je váš další příkaz?} \\
            }
        }
    }
    & \\ \\

    \theutterance \stepcounter{utterance}  
    & \multicolumn{4}{p{0.6\linewidth}}{
        \cellcolor[rgb]{0.9,0.9,0.9}{
            \makecell[{{p{\linewidth}}}]{
                \texttt{\tiny{[P1$\rangle$GM]}}
                \texttt{ŘEKNI: To je skvělý a velmi přehledný plán! Souhlasím s konfigurací v řádku 2: A(1,2), M(2,2), T(3,2), S(4,2), P(5,2), O(6,2), D(7,2). Začnu tedy posouvat své objekty na tyto pozice. Nejdříve se pokusím o objekt A.} \\
            }
        }
    }
    & & \\ \\

    \theutterance \stepcounter{utterance}  
    & & & \multicolumn{2}{p{0.3\linewidth}}{
        \cellcolor[rgb]{0.9,0.9,0.9}{
            \makecell[{{p{\linewidth}}}]{
                \texttt{\tiny{[GM$|$GM]}}
                \texttt{ŘEKNI: To je skvělý a velmi přehledný plán! Souhlasím s konfigurací v řádku 2: A(1,2), M(2,2), T(3,2), S(4,2), P(5,2), O(6,2), D(7,2). Začnu tedy posouvat své objekty na tyto pozice. Nejdříve se pokusím o objekt A.} \\
            }
        }
    }
    & & \\ \\

    \theutterance \stepcounter{utterance}  
    & & \multicolumn{4}{p{0.6\linewidth}}{
        \cellcolor[rgb]{0.9,0.9,0.9}{
            \makecell[{{p{\linewidth}}}]{
                \texttt{\tiny{[GM$\rangle$P2]}}
                \texttt{Vaše zpráva byla předána druhému hráči.} \\
\\ 
\\ 
\texttt{Právě hrajete kolo 2 z maximálně 28.} \\
\texttt{Dohromady jste nasbírali 0 z 16 pokut.} \\
\texttt{Druhý hráč poslal tuto zprávu:} \\
\\ 
\\ 
\texttt{"To je skvělý a velmi přehledný plán! Souhlasím s konfigurací v řádku 2: A(1,2), M(2,2), T(3,2), S(4,2), P(5,2), O(6,2), D(7,2). Začnu tedy posouvat své objekty na tyto pozice. Nejdříve se pokusím o objekt A."} \\
\\ 
\\ 
\texttt{Jaký je váš další příkaz?} \\
            }
        }
    }
    & \\ \\

    \theutterance \stepcounter{utterance}  
    & & & \multicolumn{4}{p{0.6\linewidth}}{
        \cellcolor[rgb]{0.9,0.9,0.9}{
            \makecell[{{p{\linewidth}}}]{
                \texttt{\tiny{[GM$\langle$P2]}}
                \texttt{POHYB: A, (1, 2)} \\
            }
        }
    }
     \\ \\

    \theutterance \stepcounter{utterance}  
    & & & \multicolumn{2}{p{0.3\linewidth}}{
        \cellcolor[rgb]{0.9,0.9,0.9}{
            \makecell[{{p{\linewidth}}}]{
                \texttt{\tiny{[GM$|$GM]}}
                \texttt{POHYB: A, (1, 2)} \\
            }
        }
    }
    & & \\ \\

    \theutterance \stepcounter{utterance}  
    & & & \multicolumn{2}{p{0.3\linewidth}}{
        \cellcolor[rgb]{0.9,0.9,0.9}{
            \makecell[{{p{\linewidth}}}]{
                \texttt{\tiny{[GM$|$GM]}}
                \texttt{Penalizace: (1,2) není prázdný, ale obsahuje '─'.} \\
\\ 
\texttt{Dohromady jste nasbírali 1 z 16 pokut.} \\
\texttt{Zkuste to prosím znovu!} \\
            }
        }
    }
    & & \\ \\

    \theutterance \stepcounter{utterance}  
    & & \multicolumn{4}{p{0.6\linewidth}}{
        \cellcolor[rgb]{0.9,0.9,0.9}{
            \makecell[{{p{\linewidth}}}]{
                \texttt{\tiny{[GM$\rangle$P2]}}
                \texttt{Penalizace: (1,2) není prázdný, ale obsahuje '─'.} \\
\\ 
\texttt{Dohromady jste nasbírali 1 z 16 pokut.} \\
\texttt{Zkuste to prosím znovu!} \\
            }
        }
    }
    & \\ \\

    \theutterance \stepcounter{utterance}  
    & & & \multicolumn{4}{p{0.6\linewidth}}{
        \cellcolor[rgb]{0.9,0.9,0.9}{
            \makecell[{{p{\linewidth}}}]{
                \texttt{\tiny{[GM$\langle$P2]}}
                \texttt{POHYB: A, (1, 3)} \\
            }
        }
    }
     \\ \\

    \theutterance \stepcounter{utterance}  
    & & & \multicolumn{2}{p{0.3\linewidth}}{
        \cellcolor[rgb]{0.9,0.9,0.9}{
            \makecell[{{p{\linewidth}}}]{
                \texttt{\tiny{[GM$|$GM]}}
                \texttt{POHYB: A, (1, 3)} \\
            }
        }
    }
    & & \\ \\

    \theutterance \stepcounter{utterance}  
    & & & \multicolumn{2}{p{0.3\linewidth}}{
        \cellcolor[rgb]{0.9,0.9,0.9}{
            \makecell[{{p{\linewidth}}}]{
                \texttt{\tiny{[GM$|$GM]}}
                \texttt{Penalizace: (1,3) není prázdný, ale obsahuje 'A'.} \\
\\ 
\texttt{Dohromady jste nasbírali 2 z 16 pokut.} \\
\texttt{Zkuste to prosím znovu!} \\
            }
        }
    }
    & & \\ \\

    \theutterance \stepcounter{utterance}  
    & & \multicolumn{4}{p{0.6\linewidth}}{
        \cellcolor[rgb]{0.9,0.9,0.9}{
            \makecell[{{p{\linewidth}}}]{
                \texttt{\tiny{[GM$\rangle$P2]}}
                \texttt{Penalizace: (1,3) není prázdný, ale obsahuje 'A'.} \\
\\ 
\texttt{Dohromady jste nasbírali 2 z 16 pokut.} \\
\texttt{Zkuste to prosím znovu!} \\
            }
        }
    }
    & \\ \\

    \theutterance \stepcounter{utterance}  
    & & & \multicolumn{4}{p{0.6\linewidth}}{
        \cellcolor[rgb]{0.9,0.9,0.9}{
            \makecell[{{p{\linewidth}}}]{
                \texttt{\tiny{[GM$\langle$P2]}}
                \texttt{POHYB: A, (1, 2)} \\
\\ 
\texttt{Wait, I see the mistake. The cell (1,2) contains a border line `╟`. I must choose a cell that is truly empty (`◌`). Looking at my grid:} \\
\texttt{(1,1) is `║`} \\
\texttt{(1,2) is `╟`} \\
\texttt{(1,3) is `║`} \\
\texttt{(1,4) is `╟`} \\
\texttt{(1,5) is `║`} \\
\texttt{(1,6) is `║`} \\
\texttt{(1,7) is `╚`} \\
\\ 
\texttt{Let me re{-}examine the grid for empty cells in column 1:} \\
\texttt{(1,1) is `║`} \\
\texttt{(1,2) is `╟`} \\
\texttt{(1,3) is `║`} \\
\texttt{(1,4) is `╟`} \\
\texttt{(1,5) is `║`} \\
\texttt{(1,6) is `║`} \\
\texttt{(1,7) is `╚`} \\
\\ 
\texttt{Actually, looking at the grid provided:} \\
\texttt{Row 1: `║◌│◌T◌│◌║` {-}> (1,1) is `◌`, (2,1) is `│`, (3,1) is `◌`, (4,1) is `T`, (5,1) is `◌`, (6,1) is `│`, (7,1) is `◌`} \\
\texttt{Row 2: `╟─┼──┬┘◌║` {-}> (1,2) is `╟`, (2,2) is `─`, (3,2) is `┼`, (4,2) is `─`, (5,2) is `─`, (6,2) is `┬`, (7,2) is `┘` (Wait, the grid is 7 wide, but the characters are tricky).} \\
\\ 
\texttt{Let's re{-}read the grid carefully:} \\
\texttt{Row 1: `║ 1 2 3 4 5 6 7 ║` (Wait, the numbers are labels)} \\
\texttt{`║◌│◌T◌│◌║ 1`} \\
\texttt{Col 1: `◌`, Col 2: `│`, Col 3: `◌`, Col 4: `T`, Col 5: `◌`, Col 6: `│`, Col 7: `◌`} \\
\\ 
\texttt{Row 2: `╟─┼──┬┘◌║ 2`} \\
\texttt{Col 1: `╟`, Col 2: `─`, Col 3: `┼`, Col 4: `─`, Col 5: `─`, Col 6: `┬`, Col 7: `┘` (Wait, the last char is `◌`? No, the last char is `║`).} \\
\texttt{Let's look at the right edge: `║ 1`, `║ 2`, `║ 3`, `║ 4`, `║ 5`, `║ 6`, `║ 7`.} \\
\texttt{The `║` is the border. The column 7 is the last column before the border.} \\
\\ 
\texttt{Let's re{-}map my grid:} \\
\texttt{(1,1)=◌, (2,1)=│, (3,1)=◌, (4,1)=T, (5,1)=◌, (6,1)=│, (7,1)=◌} \\
\texttt{(1,2)=╟, (2,2)=─, (3,2)=┼, (4,2)=─, (5,2)=─, (6,2)=┬, (7,2)=◌ (Wait, the `┘` is part of the border structure)} \\
\texttt{(1,3)=║, (2,3)=A, (3,3)=│, (4,3)=M, (5,3)=◌, (6,3)=│, (7,3)=◌} \\
\texttt{(1,4)=╟, (2,4)=─, (3,4)=┴, (4,4)=─, (5,4)=─, (6,4)=┘, (7,4)=S} \\
\texttt{(1,5)=║, (2,5)=◌, (3,5)=◌, (4,5)=◌, (5,5)=◌, (6,5)=◌, (7,5)=◌} \\
\texttt{(1,6)=║, (2,6)=◌, (3,6)=P, (4,6)=◌, (5,6)=◌, (6,6)=◌, (7,6)=◌} \\
\texttt{(1,7)=║, (2,7)=◌, (3,7)=◌, (4,7)=O, (5,7)=◌, (6,7)=D, (7,7)=◌} \\
\\ 
\texttt{My target for A was (1,2). But (1,2) is `╟`.} \\
\texttt{My target for M was (2,2). But (2,2) is `─`.} \\
\texttt{My target for T was (3,2). But (3,2) is `┼`.} \\
\\ 
\texttt{The user's grid is very tight. I need to find empty `◌` cells.} \\
\texttt{Empty cells in my grid:} \\
\texttt{(1,1), (3,1), (5,1), (7,1)} \\
\texttt{(7,2)} \\
\texttt{(5,3), (7,3)} \\
\texttt{(7,4) is S} \\
\texttt{(2,5), (3,5), (4,5), (5,5), (6,5), (7,5)} \\
\texttt{(2,6), (4,6), (5,6), (6,6), (7,6)} \\
\texttt{(2,7), (3,7), (5,7), (7,7)} \\
\\ 
\texttt{I need to pick a new target configuration that only uses `◌` cells.} \\
\texttt{Let's try Row 5 for everything.} \\
\texttt{A(1,5), M(2,5), T(3,5), S(4,5), P(5,5), O(6,5), D(7,5).} \\
\texttt{Wait, I must communicate this change. But I only have one command. I'll try to move A to a valid empty cell first.} \\
\\ 
\texttt{Let's try to move A to (2,5).} \\
\texttt{Wait, I'll try to move A to (1,1) first to get it out of the way, or just pick a valid empty cell.} \\
\texttt{Actually, I'll try to move A to (2,5).} \\
\\ 
\texttt{POHYB: A, (2, 5)} \\
            }
        }
    }
     \\ \\

    \theutterance \stepcounter{utterance}  
    & & & \multicolumn{2}{p{0.3\linewidth}}{
        \cellcolor[rgb]{0.9,0.9,0.9}{
            \makecell[{{p{\linewidth}}}]{
                \texttt{\tiny{[GM$|$GM]}}
                \texttt{POHYB: A, (1, 2)} \\
\\ 
\texttt{Wait, I see the mistake. The cell (1,2) contains a border line `╟`. I must choose a cell that is truly empty (`◌`). Looking at my grid:} \\
\texttt{(1,1) is `║`} \\
\texttt{(1,2) is `╟`} \\
\texttt{(1,3) is `║`} \\
\texttt{(1,4) is `╟`} \\
\texttt{(1,5) is `║`} \\
\texttt{(1,6) is `║`} \\
\texttt{(1,7) is `╚`} \\
\\ 
\texttt{Let me re{-}examine the grid for empty cells in column 1:} \\
\texttt{(1,1) is `║`} \\
\texttt{(1,2) is `╟`} \\
\texttt{(1,3) is `║`} \\
\texttt{(1,4) is `╟`} \\
\texttt{(1,5) is `║`} \\
\texttt{(1,6) is `║`} \\
\texttt{(1,7) is `╚`} \\
\\ 
\texttt{Actually, looking at the grid provided:} \\
\texttt{Row 1: `║◌│◌T◌│◌║` {-}> (1,1) is `◌`, (2,1) is `│`, (3,1) is `◌`, (4,1) is `T`, (5,1) is `◌`, (6,1) is `│`, (7,1) is `◌`} \\
\texttt{Row 2: `╟─┼──┬┘◌║` {-}> (1,2) is `╟`, (2,2) is `─`, (3,2) is `┼`, (4,2) is `─`, (5,2) is `─`, (6,2) is `┬`, (7,2) is `┘` (Wait, the grid is 7 wide, but the characters are tricky).} \\
\\ 
\texttt{Let's re{-}read the grid carefully:} \\
\texttt{Row 1: `║ 1 2 3 4 5 6 7 ║` (Wait, the numbers are labels)} \\
\texttt{`║◌│◌T◌│◌║ 1`} \\
\texttt{Col 1: `◌`, Col 2: `│`, Col 3: `◌`, Col 4: `T`, Col 5: `◌`, Col 6: `│`, Col 7: `◌`} \\
\\ 
\texttt{Row 2: `╟─┼──┬┘◌║ 2`} \\
\texttt{Col 1: `╟`, Col 2: `─`, Col 3: `┼`, Col 4: `─`, Col 5: `─`, Col 6: `┬`, Col 7: `┘` (Wait, the last char is `◌`? No, the last char is `║`).} \\
\texttt{Let's look at the right edge: `║ 1`, `║ 2`, `║ 3`, `║ 4`, `║ 5`, `║ 6`, `║ 7`.} \\
\texttt{The `║` is the border. The column 7 is the last column before the border.} \\
\\ 
\texttt{Let's re{-}map my grid:} \\
\texttt{(1,1)=◌, (2,1)=│, (3,1)=◌, (4,1)=T, (5,1)=◌, (6,1)=│, (7,1)=◌} \\
\texttt{(1,2)=╟, (2,2)=─, (3,2)=┼, (4,2)=─, (5,2)=─, (6,2)=┬, (7,2)=◌ (Wait, the `┘` is part of the border structure)} \\
\texttt{(1,3)=║, (2,3)=A, (3,3)=│, (4,3)=M, (5,3)=◌, (6,3)=│, (7,3)=◌} \\
\texttt{(1,4)=╟, (2,4)=─, (3,4)=┴, (4,4)=─, (5,4)=─, (6,4)=┘, (7,4)=S} \\
\texttt{(1,5)=║, (2,5)=◌, (3,5)=◌, (4,5)=◌, (5,5)=◌, (6,5)=◌, (7,5)=◌} \\
\texttt{(1,6)=║, (2,6)=◌, (3,6)=P, (4,6)=◌, (5,6)=◌, (6,6)=◌, (7,6)=◌} \\
\texttt{(1,7)=║, (2,7)=◌, (3,7)=◌, (4,7)=O, (5,7)=◌, (6,7)=D, (7,7)=◌} \\
\\ 
\texttt{My target for A was (1,2). But (1,2) is `╟`.} \\
\texttt{My target for M was (2,2). But (2,2) is `─`.} \\
\texttt{My target for T was (3,2). But (3,2) is `┼`.} \\
\\ 
\texttt{The user's grid is very tight. I need to find empty `◌` cells.} \\
\texttt{Empty cells in my grid:} \\
\texttt{(1,1), (3,1), (5,1), (7,1)} \\
\texttt{(7,2)} \\
\texttt{(5,3), (7,3)} \\
\texttt{(7,4) is S} \\
\texttt{(2,5), (3,5), (4,5), (5,5), (6,5), (7,5)} \\
\texttt{(2,6), (4,6), (5,6), (6,6), (7,6)} \\
\texttt{(2,7), (3,7), (5,7), (7,7)} \\
\\ 
\texttt{I need to pick a new target configuration that only uses `◌` cells.} \\
\texttt{Let's try Row 5 for everything.} \\
\texttt{A(1,5), M(2,5), T(3,5), S(4,5), P(5,5), O(6,5), D(7,5).} \\
\texttt{Wait, I must communicate this change. But I only have one command. I'll try to move A to a valid empty cell first.} \\
\\ 
\texttt{Let's try to move A to (2,5).} \\
\texttt{Wait, I'll try to move A to (1,1) first to get it out of the way, or just pick a valid empty cell.} \\
\texttt{Actually, I'll try to move A to (2,5).} \\
\\ 
\texttt{POHYB: A, (2, 5)} \\
            }
        }
    }
    & & \\ \\

    \theutterance \stepcounter{utterance}  
    & & & \multicolumn{2}{p{0.3\linewidth}}{
        \cellcolor[rgb]{0.9,0.9,0.9}{
            \makecell[{{p{\linewidth}}}]{
                \texttt{\tiny{[GM$|$GM]}}
                \texttt{Invalid format: head is not empty: 'POHYB: A, (1, 2)} \\
\\ 
\texttt{Wait, I see the mistake. The cell (1,2) contains a border line ╟. I must choose a cell that is truly empty (◌). Looking at my grid:} \\
\texttt{(1,1) is ║} \\
\texttt{(1,2) is ╟} \\
\texttt{(1,3) is ║} \\
\texttt{(1,4) is ╟} \\
\texttt{(1,5) is ║} \\
\texttt{(1,6) is ║} \\
\texttt{(1,7) is ╚} \\
\\ 
\texttt{Let me re{-}examine the grid for empty cells in column 1:} \\
\texttt{(1,1) is ║} \\
\texttt{(1,2) is ╟} \\
\texttt{(1,3) is ║} \\
\texttt{(1,4) is ╟} \\
\texttt{(1,5) is ║} \\
\texttt{(1,6) is ║} \\
\texttt{(1,7) is ╚} \\
\\ 
\texttt{Actually, looking at the grid provided:} \\
\texttt{Row 1: ║◌│◌T◌│◌║ {-}> (1,1) is ◌, (2,1) is │, (3,1) is ◌, (4,1) is T, (5,1) is ◌, (6,1) is │, (7,1) is ◌} \\
\texttt{Row 2: ╟─┼──┬┘◌║ {-}> (1,2) is ╟, (2,2) is ─, (3,2) is ┼, (4,2) is ─, (5,2) is ─, (6,2) is ┬, (7,2) is ┘ (Wait, the grid is 7 wide, but the characters are tricky).} \\
\\ 
\texttt{Let's re{-}read the grid carefully:} \\
\texttt{Row 1: ║ 1 2 3 4 5 6 7 ║ (Wait, the numbers are labels)} \\
\texttt{║◌│◌T◌│◌║ 1} \\
\texttt{Col 1: ◌, Col 2: │, Col 3: ◌, Col 4: T, Col 5: ◌, Col 6: │, Col 7: ◌} \\
\\ 
\texttt{Row 2: ╟─┼──┬┘◌║ 2} \\
\texttt{Col 1: ╟, Col 2: ─, Col 3: ┼, Col 4: ─, Col 5: ─, Col 6: ┬, Col 7: ┘ (Wait, the last char is ◌? No, the last char is ║).} \\
\texttt{Let's look at the right edge: ║ 1, ║ 2, ║ 3, ║ 4, ║ 5, ║ 6, ║ 7.} \\
\texttt{The ║ is the border. The column 7 is the last column before the border.} \\
\\ 
\texttt{Let's re{-}map my grid:} \\
\texttt{(1,1)=◌, (2,1)=│, (3,1)=◌, (4,1)=T, (5,1)=◌, (6,1)=│, (7,1)=◌} \\
\texttt{(1,2)=╟, (2,2)=─, (3,2)=┼, (4,2)=─, (5,2)=─, (6,2)=┬, (7,2)=◌ (Wait, the ┘ is part of the border structure)} \\
\texttt{(1,3)=║, (2,3)=A, (3,3)=│, (4,3)=M, (5,3)=◌, (6,3)=│, (7,3)=◌} \\
\texttt{(1,4)=╟, (2,4)=─, (3,4)=┴, (4,4)=─, (5,4)=─, (6,4)=┘, (7,4)=S} \\
\texttt{(1,5)=║, (2,5)=◌, (3,5)=◌, (4,5)=◌, (5,5)=◌, (6,5)=◌, (7,5)=◌} \\
\texttt{(1,6)=║, (2,6)=◌, (3,6)=P, (4,6)=◌, (5,6)=◌, (6,6)=◌, (7,6)=◌} \\
\texttt{(1,7)=║, (2,7)=◌, (3,7)=◌, (4,7)=O, (5,7)=◌, (6,7)=D, (7,7)=◌} \\
\\ 
\texttt{My target for A was (1,2). But (1,2) is ╟.} \\
\texttt{My target for M was (2,2). But (2,2) is ─.} \\
\texttt{My target for T was (3,2). But (3,2) is ┼.} \\
\\ 
\texttt{The user's grid is very tight. I need to find empty ◌ cells.} \\
\texttt{Empty cells in my grid:} \\
\texttt{(1,1), (3,1), (5,1), (7,1)} \\
\texttt{(7,2)} \\
\texttt{(5,3), (7,3)} \\
\texttt{(7,4) is S} \\
\texttt{(2,5), (3,5), (4,5), (5,5), (6,5), (7,5)} \\
\texttt{(2,6), (4,6), (5,6), (6,6), (7,6)} \\
\texttt{(2,7), (3,7), (5,7), (7,7)} \\
\\ 
\texttt{I need to pick a new target configuration that only uses ◌ cells.} \\
\texttt{Let's try Row 5 for everything.} \\
\texttt{A(1,5), M(2,5), T(3,5), S(4,5), P(5,5), O(6,5), D(7,5).} \\
\texttt{Wait, I must communicate this change. But I only have one command. I'll try to move A to a valid empty cell first.} \\
\\ 
\texttt{Let's try to move A to (2,5).} \\
\texttt{Wait, I'll try to move A to (1,1) first to get it out of the way, or just pick a valid empty cell.} \\
\texttt{Actually, I'll try to move A to (2,5).} \\
\\ 
\texttt{'} \\
            }
        }
    }
    & & \\ \\

    \theutterance \stepcounter{utterance}  
    & & \multicolumn{4}{p{0.6\linewidth}}{
        \cellcolor[rgb]{0.9,0.9,0.9}{
            \makecell[{{p{\linewidth}}}]{
                \texttt{\tiny{[GM$\rangle$P2]}}
                \texttt{Trest: Vaše zpráva nesmí obsahovat nic před příkazem!} \\
\\ 
\texttt{Ujistěte se, že vaše odpověď obsahuje pouze `ŘEKNI: <ZPRÁVA>` nebo `POHYB: <OBJEKT>, (<X>, <Y>)` a nic jiného!} \\
\\ 
\texttt{Dohromady jste nasbírali 3 z 16 pokut.} \\
\texttt{Zkuste to prosím znovu!} \\
            }
        }
    }
    & \\ \\

    \theutterance \stepcounter{utterance}  
    & & & \multicolumn{4}{p{0.6\linewidth}}{
        \cellcolor[rgb]{0.9,0.9,0.9}{
            \makecell[{{p{\linewidth}}}]{
                \texttt{\tiny{[GM$\langle$P2]}}
                \texttt{POHYB: A, (2, 5)} \\
            }
        }
    }
     \\ \\

    \theutterance \stepcounter{utterance}  
    & & & \multicolumn{2}{p{0.3\linewidth}}{
        \cellcolor[rgb]{0.9,0.9,0.9}{
            \makecell[{{p{\linewidth}}}]{
                \texttt{\tiny{[GM$|$GM]}}
                \texttt{POHYB: A, (2, 5)} \\
            }
        }
    }
    & & \\ \\

    \theutterance \stepcounter{utterance}  
    & & & \multicolumn{2}{p{0.3\linewidth}}{
        \cellcolor[rgb]{0.9,0.9,0.9}{
            \makecell[{{p{\linewidth}}}]{
                \texttt{\tiny{[GM$|$GM]}}
                \texttt{A moved to (2, 5)} \\
            }
        }
    }
    & & \\ \\

    \theutterance \stepcounter{utterance}  
    & & \multicolumn{4}{p{0.6\linewidth}}{
        \cellcolor[rgb]{0.9,0.9,0.9}{
            \makecell[{{p{\linewidth}}}]{
                \texttt{\tiny{[P1$\langle$GM]}}
                \texttt{Vaše zpráva byla předána druhému hráči.} \\
\\ 
\\ 
\texttt{Právě hrajete kolo 2 z maximálně 28.} \\
\texttt{Dohromady jste nasbírali 3 z 16 pokut.} \\
\texttt{Druhý hráč přesunul objekt na své mřížce.} \\
\\ 
\\ 
\texttt{Jaký je váš další příkaz?} \\
            }
        }
    }
    & \\ \\

    \theutterance \stepcounter{utterance}  
    & \multicolumn{4}{p{0.6\linewidth}}{
        \cellcolor[rgb]{0.9,0.9,0.9}{
            \makecell[{{p{\linewidth}}}]{
                \texttt{\tiny{[P1$\rangle$GM]}}
                \texttt{POHYB: A, (1, 2)} \\
            }
        }
    }
    & & \\ \\

    \theutterance \stepcounter{utterance}  
    & & & \multicolumn{2}{p{0.3\linewidth}}{
        \cellcolor[rgb]{0.9,0.9,0.9}{
            \makecell[{{p{\linewidth}}}]{
                \texttt{\tiny{[GM$|$GM]}}
                \texttt{POHYB: A, (1, 2)} \\
            }
        }
    }
    & & \\ \\

    \theutterance \stepcounter{utterance}  
    & & & \multicolumn{2}{p{0.3\linewidth}}{
        \cellcolor[rgb]{0.9,0.9,0.9}{
            \makecell[{{p{\linewidth}}}]{
                \texttt{\tiny{[GM$|$GM]}}
                \texttt{Penalizace: (1,2) není prázdný, ale obsahuje '─'.} \\
\\ 
\texttt{Dohromady jste nasbírali 4 z 16 pokut.} \\
\texttt{Zkuste to prosím znovu!} \\
            }
        }
    }
    & & \\ \\

    \theutterance \stepcounter{utterance}  
    & & \multicolumn{4}{p{0.6\linewidth}}{
        \cellcolor[rgb]{0.9,0.9,0.9}{
            \makecell[{{p{\linewidth}}}]{
                \texttt{\tiny{[P1$\langle$GM]}}
                \texttt{Penalizace: (1,2) není prázdný, ale obsahuje '─'.} \\
\\ 
\texttt{Dohromady jste nasbírali 4 z 16 pokut.} \\
\texttt{Zkuste to prosím znovu!} \\
            }
        }
    }
    & \\ \\

    \theutterance \stepcounter{utterance}  
    & \multicolumn{4}{p{0.6\linewidth}}{
        \cellcolor[rgb]{0.9,0.9,0.9}{
            \makecell[{{p{\linewidth}}}]{
                \texttt{\tiny{[P1$\rangle$GM]}}
                \texttt{POHYB: A, (1, 3)} \\
            }
        }
    }
    & & \\ \\

    \theutterance \stepcounter{utterance}  
    & & & \multicolumn{2}{p{0.3\linewidth}}{
        \cellcolor[rgb]{0.9,0.9,0.9}{
            \makecell[{{p{\linewidth}}}]{
                \texttt{\tiny{[GM$|$GM]}}
                \texttt{POHYB: A, (1, 3)} \\
            }
        }
    }
    & & \\ \\

    \theutterance \stepcounter{utterance}  
    & & & \multicolumn{2}{p{0.3\linewidth}}{
        \cellcolor[rgb]{0.9,0.9,0.9}{
            \makecell[{{p{\linewidth}}}]{
                \texttt{\tiny{[GM$|$GM]}}
                \texttt{A moved to (1, 3)} \\
            }
        }
    }
    & & \\ \\

    \theutterance \stepcounter{utterance}  
    & & \multicolumn{4}{p{0.6\linewidth}}{
        \cellcolor[rgb]{0.9,0.9,0.9}{
            \makecell[{{p{\linewidth}}}]{
                \texttt{\tiny{[GM$\rangle$P2]}}
                \texttt{Úspěšně přesunuto 'A' do (2,5). Vaše aktualizovaná mřížka vypadá takto:} \\
\texttt{```} \\
\texttt{ 1234567} \\
\texttt{╔═╤═══╤═╗} \\
\texttt{║◌│◌T◌│◌║ 1} \\
\texttt{╟─┼──┬┘◌║ 2} \\
\texttt{║◌│M◌│◌◌║ 3} \\
\texttt{╟─┴──┘◌S║ 4} \\
\texttt{║◌A◌◌◌◌◌║ 5} \\
\texttt{║◌P◌◌◌◌◌║ 6} \\
\texttt{║◌◌O◌D◌◌║ 7} \\
\texttt{╚═══════╝} \\
\texttt{```} \\
\texttt{Právě hrajete kolo 3 z maximálně 28.} \\
\texttt{Dohromady jste nasbírali 4 z 16 pokut.} \\
\texttt{Druhý hráč přesunul objekt na své mřížce.} \\
\\ 
\\ 
\texttt{Jaký je váš další příkaz?} \\
            }
        }
    }
    & \\ \\

    \theutterance \stepcounter{utterance}  
    & & & \multicolumn{4}{p{0.6\linewidth}}{
        \cellcolor[rgb]{0.9,0.9,0.9}{
            \makecell[{{p{\linewidth}}}]{
                \texttt{\tiny{[GM$\langle$P2]}}
                \texttt{POHYB: M, (3, 5)} \\
            }
        }
    }
     \\ \\

    \theutterance \stepcounter{utterance}  
    & & & \multicolumn{2}{p{0.3\linewidth}}{
        \cellcolor[rgb]{0.9,0.9,0.9}{
            \makecell[{{p{\linewidth}}}]{
                \texttt{\tiny{[GM$|$GM]}}
                \texttt{POHYB: M, (3, 5)} \\
            }
        }
    }
    & & \\ \\

    \theutterance \stepcounter{utterance}  
    & & & \multicolumn{2}{p{0.3\linewidth}}{
        \cellcolor[rgb]{0.9,0.9,0.9}{
            \makecell[{{p{\linewidth}}}]{
                \texttt{\tiny{[GM$|$GM]}}
                \texttt{M moved to (3, 5)} \\
            }
        }
    }
    & & \\ \\

    \theutterance \stepcounter{utterance}  
    & & \multicolumn{4}{p{0.6\linewidth}}{
        \cellcolor[rgb]{0.9,0.9,0.9}{
            \makecell[{{p{\linewidth}}}]{
                \texttt{\tiny{[P1$\langle$GM]}}
                \texttt{Úspěšně přesunuto 'A' do (1,3). Vaše aktualizovaná mřížka vypadá takto:} \\
\texttt{```} \\
\texttt{ 1234567} \\
\texttt{╔═╤═══╤═╗} \\
\texttt{║◌│◌DT│◌║ 1} \\
\texttt{╟─┼──┬┘◌║ 2} \\
\texttt{║A│◌◌│◌◌║ 3} \\
\texttt{╟─┴──┘◌O║ 4} \\
\texttt{║M◌◌P◌◌◌║ 5} \\
\texttt{║◌◌◌◌◌◌◌║ 6} \\
\texttt{║S◌◌◌◌◌◌║ 7} \\
\texttt{╚═══════╝} \\
\texttt{```} \\
\texttt{Právě hrajete kolo 3 z maximálně 28.} \\
\texttt{Dohromady jste nasbírali 4 z 16 pokut.} \\
\texttt{Druhý hráč přesunul objekt na své mřížce.} \\
\\ 
\\ 
\texttt{Jaký je váš další příkaz?} \\
            }
        }
    }
    & \\ \\

    \theutterance \stepcounter{utterance}  
    & \multicolumn{4}{p{0.6\linewidth}}{
        \cellcolor[rgb]{0.9,0.9,0.9}{
            \makecell[{{p{\linewidth}}}]{
                \texttt{\tiny{[P1$\rangle$GM]}}
                \texttt{POHYB: M, (2, 2)} \\
            }
        }
    }
    & & \\ \\

    \theutterance \stepcounter{utterance}  
    & & & \multicolumn{2}{p{0.3\linewidth}}{
        \cellcolor[rgb]{0.9,0.9,0.9}{
            \makecell[{{p{\linewidth}}}]{
                \texttt{\tiny{[GM$|$GM]}}
                \texttt{POHYB: M, (2, 2)} \\
            }
        }
    }
    & & \\ \\

    \theutterance \stepcounter{utterance}  
    & & & \multicolumn{2}{p{0.3\linewidth}}{
        \cellcolor[rgb]{0.9,0.9,0.9}{
            \makecell[{{p{\linewidth}}}]{
                \texttt{\tiny{[GM$|$GM]}}
                \texttt{Penalizace: (2,2) není prázdný, ale obsahuje '┼'.} \\
\\ 
\texttt{Dohromady jste nasbírali 5 z 16 pokut.} \\
\texttt{Zkuste to prosím znovu!} \\
            }
        }
    }
    & & \\ \\

    \theutterance \stepcounter{utterance}  
    & & \multicolumn{4}{p{0.6\linewidth}}{
        \cellcolor[rgb]{0.9,0.9,0.9}{
            \makecell[{{p{\linewidth}}}]{
                \texttt{\tiny{[P1$\langle$GM]}}
                \texttt{Penalizace: (2,2) není prázdný, ale obsahuje '┼'.} \\
\\ 
\texttt{Dohromady jste nasbírali 5 z 16 pokut.} \\
\texttt{Zkuste to prosím znovu!} \\
            }
        }
    }
    & \\ \\

    \theutterance \stepcounter{utterance}  
    & \multicolumn{4}{p{0.6\linewidth}}{
        \cellcolor[rgb]{0.9,0.9,0.9}{
            \makecell[{{p{\linewidth}}}]{
                \texttt{\tiny{[P1$\rangle$GM]}}
                \texttt{POHYB: M, (1, 5)} \\
            }
        }
    }
    & & \\ \\

    \theutterance \stepcounter{utterance}  
    & & & \multicolumn{2}{p{0.3\linewidth}}{
        \cellcolor[rgb]{0.9,0.9,0.9}{
            \makecell[{{p{\linewidth}}}]{
                \texttt{\tiny{[GM$|$GM]}}
                \texttt{POHYB: M, (1, 5)} \\
            }
        }
    }
    & & \\ \\

    \theutterance \stepcounter{utterance}  
    & & & \multicolumn{2}{p{0.3\linewidth}}{
        \cellcolor[rgb]{0.9,0.9,0.9}{
            \makecell[{{p{\linewidth}}}]{
                \texttt{\tiny{[GM$|$GM]}}
                \texttt{Penalizace: (1,5) není prázdný, ale obsahuje 'M'.} \\
\\ 
\texttt{Dohromady jste nasbírali 6 z 16 pokut.} \\
\texttt{Zkuste to prosím znovu!} \\
            }
        }
    }
    & & \\ \\

    \theutterance \stepcounter{utterance}  
    & & \multicolumn{4}{p{0.6\linewidth}}{
        \cellcolor[rgb]{0.9,0.9,0.9}{
            \makecell[{{p{\linewidth}}}]{
                \texttt{\tiny{[P1$\langle$GM]}}
                \texttt{Penalizace: (1,5) není prázdný, ale obsahuje 'M'.} \\
\\ 
\texttt{Dohromady jste nasbírali 6 z 16 pokut.} \\
\texttt{Zkuste to prosím znovu!} \\
            }
        }
    }
    & \\ \\

    \theutterance \stepcounter{utterance}  
    & \multicolumn{4}{p{0.6\linewidth}}{
        \cellcolor[rgb]{0.9,0.9,0.9}{
            \makecell[{{p{\linewidth}}}]{
                \texttt{\tiny{[P1$\rangle$GM]}}
                \texttt{POHYB: M, (2, 5)} \\
            }
        }
    }
    & & \\ \\

    \theutterance \stepcounter{utterance}  
    & & & \multicolumn{2}{p{0.3\linewidth}}{
        \cellcolor[rgb]{0.9,0.9,0.9}{
            \makecell[{{p{\linewidth}}}]{
                \texttt{\tiny{[GM$|$GM]}}
                \texttt{POHYB: M, (2, 5)} \\
            }
        }
    }
    & & \\ \\

    \theutterance \stepcounter{utterance}  
    & & & \multicolumn{2}{p{0.3\linewidth}}{
        \cellcolor[rgb]{0.9,0.9,0.9}{
            \makecell[{{p{\linewidth}}}]{
                \texttt{\tiny{[GM$|$GM]}}
                \texttt{M moved to (2, 5)} \\
            }
        }
    }
    & & \\ \\

    \theutterance \stepcounter{utterance}  
    & & \multicolumn{4}{p{0.6\linewidth}}{
        \cellcolor[rgb]{0.9,0.9,0.9}{
            \makecell[{{p{\linewidth}}}]{
                \texttt{\tiny{[GM$\rangle$P2]}}
                \texttt{Úspěšně přesunuto 'M' do (3,5). Vaše aktualizovaná mřížka vypadá takto:} \\
\texttt{```} \\
\texttt{ 1234567} \\
\texttt{╔═╤═══╤═╗} \\
\texttt{║◌│◌T◌│◌║ 1} \\
\texttt{╟─┼──┬┘◌║ 2} \\
\texttt{║◌│◌◌│◌◌║ 3} \\
\texttt{╟─┴──┘◌S║ 4} \\
\texttt{║◌AM◌◌◌◌║ 5} \\
\texttt{║◌P◌◌◌◌◌║ 6} \\
\texttt{║◌◌O◌D◌◌║ 7} \\
\texttt{╚═══════╝} \\
\texttt{```} \\
\texttt{Právě hrajete kolo 4 z maximálně 28.} \\
\texttt{Dohromady jste nasbírali 6 z 16 pokut.} \\
\texttt{Druhý hráč přesunul objekt na své mřížce.} \\
\\ 
\\ 
\texttt{Jaký je váš další příkaz?} \\
            }
        }
    }
    & \\ \\

    \theutterance \stepcounter{utterance}  
    & & & \multicolumn{4}{p{0.6\linewidth}}{
        \cellcolor[rgb]{0.9,0.9,0.9}{
            \makecell[{{p{\linewidth}}}]{
                \texttt{\tiny{[GM$\langle$P2]}}
                \texttt{POHYB: T, (5, 5)} \\
            }
        }
    }
     \\ \\

    \theutterance \stepcounter{utterance}  
    & & & \multicolumn{2}{p{0.3\linewidth}}{
        \cellcolor[rgb]{0.9,0.9,0.9}{
            \makecell[{{p{\linewidth}}}]{
                \texttt{\tiny{[GM$|$GM]}}
                \texttt{POHYB: T, (5, 5)} \\
            }
        }
    }
    & & \\ \\

    \theutterance \stepcounter{utterance}  
    & & & \multicolumn{2}{p{0.3\linewidth}}{
        \cellcolor[rgb]{0.9,0.9,0.9}{
            \makecell[{{p{\linewidth}}}]{
                \texttt{\tiny{[GM$|$GM]}}
                \texttt{T moved to (5, 5)} \\
            }
        }
    }
    & & \\ \\

    \theutterance \stepcounter{utterance}  
    & & \multicolumn{4}{p{0.6\linewidth}}{
        \cellcolor[rgb]{0.9,0.9,0.9}{
            \makecell[{{p{\linewidth}}}]{
                \texttt{\tiny{[P1$\langle$GM]}}
                \texttt{Úspěšně přesunuto 'M' do (2,5). Vaše aktualizovaná mřížka vypadá takto:} \\
\texttt{```} \\
\texttt{ 1234567} \\
\texttt{╔═╤═══╤═╗} \\
\texttt{║◌│◌DT│◌║ 1} \\
\texttt{╟─┼──┬┘◌║ 2} \\
\texttt{║A│◌◌│◌◌║ 3} \\
\texttt{╟─┴──┘◌O║ 4} \\
\texttt{║◌M◌P◌◌◌║ 5} \\
\texttt{║◌◌◌◌◌◌◌║ 6} \\
\texttt{║S◌◌◌◌◌◌║ 7} \\
\texttt{╚═══════╝} \\
\texttt{```} \\
\texttt{Právě hrajete kolo 4 z maximálně 28.} \\
\texttt{Dohromady jste nasbírali 6 z 16 pokut.} \\
\texttt{Druhý hráč přesunul objekt na své mřížce.} \\
\\ 
\\ 
\texttt{Jaký je váš další příkaz?} \\
            }
        }
    }
    & \\ \\

    \theutterance \stepcounter{utterance}  
    & \multicolumn{4}{p{0.6\linewidth}}{
        \cellcolor[rgb]{0.9,0.9,0.9}{
            \makecell[{{p{\linewidth}}}]{
                \texttt{\tiny{[P1$\rangle$GM]}}
                \texttt{POHYB: M, (2, 2)} \\
\\ 
\texttt{Wait, I just realized (2,2) is a junction '┼'. I must avoid the lines and junctions. Let me re{-}evaluate the target row 2.} \\
\texttt{The target row 2 is: A(1,2), M(2,2), T(3,2), S(4,2), P(5,2), O(6,2), D(7,2).} \\
\texttt{Looking at the grid, row 2 is mostly lines: `╟─┼──┬┘◌║ 2`.} \\
\texttt{The only empty cell in row 2 is (7,2).} \\
\texttt{My previous plan for row 2 was impossible because most cells in row 2 are part of the grid structure (lines/borders).} \\
\\ 
\texttt{Let's change the target row to row 3 or 6, which are completely empty.} \\
\texttt{Let's try row 6: A(1,6), M(2,6), T(3,6), S(4,6), P(5,6), O(6,6), D(7,6).} \\
\\ 
\texttt{Wait, I need to communicate this change to my partner first.} \\
\\ 
\texttt{ŘEKNI: Omlouvám se, uvědomil jsem si, že řádek 2 je plný čar a hranic, takže tam objekty neposuneme. Změňme cíl na řádek 6: A(1,6), M(2,6), T(3,6), S(4,6), P(5,6), O(6,6), D(7,6). Co na to říkáš?} \\
            }
        }
    }
    & & \\ \\

    \theutterance \stepcounter{utterance}  
    & & & \multicolumn{2}{p{0.3\linewidth}}{
        \cellcolor[rgb]{0.9,0.9,0.9}{
            \makecell[{{p{\linewidth}}}]{
                \texttt{\tiny{[GM$|$GM]}}
                \texttt{POHYB: M, (2, 2)} \\
\\ 
\texttt{Wait, I just realized (2,2) is a junction '┼'. I must avoid the lines and junctions. Let me re{-}evaluate the target row 2.} \\
\texttt{The target row 2 is: A(1,2), M(2,2), T(3,2), S(4,2), P(5,2), O(6,2), D(7,2).} \\
\texttt{Looking at the grid, row 2 is mostly lines: `╟─┼──┬┘◌║ 2`.} \\
\texttt{The only empty cell in row 2 is (7,2).} \\
\texttt{My previous plan for row 2 was impossible because most cells in row 2 are part of the grid structure (lines/borders).} \\
\\ 
\texttt{Let's change the target row to row 3 or 6, which are completely empty.} \\
\texttt{Let's try row 6: A(1,6), M(2,6), T(3,6), S(4,6), P(5,6), O(6,6), D(7,6).} \\
\\ 
\texttt{Wait, I need to communicate this change to my partner first.} \\
\\ 
\texttt{ŘEKNI: Omlouvám se, uvědomil jsem si, že řádek 2 je plný čar a hranic, takže tam objekty neposuneme. Změňme cíl na řádek 6: A(1,6), M(2,6), T(3,6), S(4,6), P(5,6), O(6,6), D(7,6). Co na to říkáš?} \\
            }
        }
    }
    & & \\ \\

    \theutterance \stepcounter{utterance}  
    & & & \multicolumn{2}{p{0.3\linewidth}}{
        \cellcolor[rgb]{0.9,0.9,0.9}{
            \makecell[{{p{\linewidth}}}]{
                \texttt{\tiny{[GM$|$GM]}}
                \texttt{Invalid response format: POHYB: M, (2, 2)} \\
\\ 
\texttt{Wait, I just realized (2,2) is a junction '┼'. I must avoid the lines and junctions. Let me re{-}evaluate the target row 2.} \\
\texttt{The target row 2 is: A(1,2), M(2,2), T(3,2), S(4,2), P(5,2), O(6,2), D(7,2).} \\
\texttt{Looking at the grid, row 2 is mostly lines: ╟─┼──┬┘◌║ 2.} \\
\texttt{The only empty cell in row 2 is (7,2).} \\
\texttt{My previous plan for row 2 was impossible because most cells in row 2 are part of the grid structure (lines/borders).} \\
\\ 
\texttt{Let's change the target row to row 3 or 6, which are completely empty.} \\
\texttt{Let's try row 6: A(1,6), M(2,6), T(3,6), S(4,6), P(5,6), O(6,6), D(7,6).} \\
\\ 
\texttt{Wait, I need to communicate this change to my partner first.} \\
\\ 
\texttt{ŘEKNI: Omlouvám se, uvědomil jsem si, že řádek 2 je plný čar a hranic, takže tam objekty neposuneme. Změňme cíl na řádek 6: A(1,6), M(2,6), T(3,6), S(4,6), P(5,6), O(6,6), D(7,6). Co na to říkáš?} \\
            }
        }
    }
    & & \\ \\

    \theutterance \stepcounter{utterance}  
    & & \multicolumn{4}{p{0.6\linewidth}}{
        \cellcolor[rgb]{0.9,0.9,0.9}{
            \makecell[{{p{\linewidth}}}]{
                \texttt{\tiny{[P1$\langle$GM]}}
                \texttt{Trest: Vaše zpráva obsahuje více než jeden příkaz!} \\
\\ 
\texttt{Ujistěte se, že vaše odpověď obsahuje pouze `ŘEKNI: <ZPRÁVA>` nebo `POHYB: <OBJEKT>, (<X>, <Y>)` a nic jiného!} \\
\\ 
\texttt{Dohromady jste nasbírali 7 z 16 pokut.} \\
\texttt{Zkuste to prosím znovu!} \\
            }
        }
    }
    & \\ \\

    \theutterance \stepcounter{utterance}  
    & \multicolumn{4}{p{0.6\linewidth}}{
        \cellcolor[rgb]{0.9,0.9,0.9}{
            \makecell[{{p{\linewidth}}}]{
                \texttt{\tiny{[P1$\rangle$GM]}}
                \texttt{ŘEKNI: Omlouvám se, uvědomil jsem si, že řádek 2 je plný čar a hranic. Změňme cíl na řádek 6: A(1,6), M(2,6), T(3,6), S(4,6), P(5,6), O(6,6), D(7,6). Co na to říkáš?} \\
            }
        }
    }
    & & \\ \\

    \theutterance \stepcounter{utterance}  
    & & & \multicolumn{2}{p{0.3\linewidth}}{
        \cellcolor[rgb]{0.9,0.9,0.9}{
            \makecell[{{p{\linewidth}}}]{
                \texttt{\tiny{[GM$|$GM]}}
                \texttt{ŘEKNI: Omlouvám se, uvědomil jsem si, že řádek 2 je plný čar a hranic. Změňme cíl na řádek 6: A(1,6), M(2,6), T(3,6), S(4,6), P(5,6), O(6,6), D(7,6). Co na to říkáš?} \\
            }
        }
    }
    & & \\ \\

    \theutterance \stepcounter{utterance}  
    & & \multicolumn{4}{p{0.6\linewidth}}{
        \cellcolor[rgb]{0.9,0.9,0.9}{
            \makecell[{{p{\linewidth}}}]{
                \texttt{\tiny{[GM$\rangle$P2]}}
                \texttt{Úspěšně přesunuto 'T' do (5,5). Vaše aktualizovaná mřížka vypadá takto:} \\
\texttt{```} \\
\texttt{ 1234567} \\
\texttt{╔═╤═══╤═╗} \\
\texttt{║◌│◌◌◌│◌║ 1} \\
\texttt{╟─┼──┬┘◌║ 2} \\
\texttt{║◌│◌◌│◌◌║ 3} \\
\texttt{╟─┴──┘◌S║ 4} \\
\texttt{║◌AM◌T◌◌║ 5} \\
\texttt{║◌P◌◌◌◌◌║ 6} \\
\texttt{║◌◌O◌D◌◌║ 7} \\
\texttt{╚═══════╝} \\
\texttt{```} \\
\texttt{Právě hrajete kolo 5 z maximálně 28.} \\
\texttt{Dohromady jste nasbírali 7 z 16 pokut.} \\
\texttt{Druhý hráč poslal tuto zprávu:} \\
\\ 
\\ 
\texttt{"Omlouvám se, uvědomil jsem si, že řádek 2 je plný čar a hranic. Změňme cíl na řádek 6: A(1,6), M(2,6), T(3,6), S(4,6), P(5,6), O(6,6), D(7,6). Co na to říkáš?"} \\
\\ 
\\ 
\texttt{Jaký je váš další příkaz?} \\
            }
        }
    }
    & \\ \\

    \theutterance \stepcounter{utterance}  
    & & & \multicolumn{4}{p{0.6\linewidth}}{
        \cellcolor[rgb]{0.9,0.9,0.9}{
            \makecell[{{p{\linewidth}}}]{
                \texttt{\tiny{[GM$\langle$P2]}}
                \texttt{ŘEKNI: To je dobrý nápad, řádek 6 je mnohem čistší. Souhlasím s novým cílem: A(1,6), M(2,6), T(3,6), S(4,6), P(5,6), O(6,6), D(7,6). Budu se podle toho řídit.} \\
            }
        }
    }
     \\ \\

    \theutterance \stepcounter{utterance}  
    & & & \multicolumn{2}{p{0.3\linewidth}}{
        \cellcolor[rgb]{0.9,0.9,0.9}{
            \makecell[{{p{\linewidth}}}]{
                \texttt{\tiny{[GM$|$GM]}}
                \texttt{ŘEKNI: To je dobrý nápad, řádek 6 je mnohem čistší. Souhlasím s novým cílem: A(1,6), M(2,6), T(3,6), S(4,6), P(5,6), O(6,6), D(7,6). Budu se podle toho řídit.} \\
            }
        }
    }
    & & \\ \\

    \theutterance \stepcounter{utterance}  
    & & \multicolumn{4}{p{0.6\linewidth}}{
        \cellcolor[rgb]{0.9,0.9,0.9}{
            \makecell[{{p{\linewidth}}}]{
                \texttt{\tiny{[P1$\langle$GM]}}
                \texttt{Vaše zpráva byla předána druhému hráči.} \\
\\ 
\\ 
\texttt{Právě hrajete kolo 5 z maximálně 28.} \\
\texttt{Dohromady jste nasbírali 7 z 16 pokut.} \\
\texttt{Druhý hráč poslal tuto zprávu:} \\
\\ 
\\ 
\texttt{"To je dobrý nápad, řádek 6 je mnohem čistší. Souhlasím s novým cílem: A(1,6), M(2,6), T(3,6), S(4,6), P(5,6), O(6,6), D(7,6). Budu se podle toho řídit."} \\
\\ 
\\ 
\texttt{Jaký je váš další příkaz?} \\
            }
        }
    }
    & \\ \\

    \theutterance \stepcounter{utterance}  
    & \multicolumn{4}{p{0.6\linewidth}}{
        \cellcolor[rgb]{0.9,0.9,0.9}{
            \makecell[{{p{\linewidth}}}]{
                \texttt{\tiny{[P1$\rangle$GM]}}
                \texttt{POHYB: A, (1, 6)} \\
            }
        }
    }
    & & \\ \\

    \theutterance \stepcounter{utterance}  
    & & & \multicolumn{2}{p{0.3\linewidth}}{
        \cellcolor[rgb]{0.9,0.9,0.9}{
            \makecell[{{p{\linewidth}}}]{
                \texttt{\tiny{[GM$|$GM]}}
                \texttt{POHYB: A, (1, 6)} \\
            }
        }
    }
    & & \\ \\

    \theutterance \stepcounter{utterance}  
    & & & \multicolumn{2}{p{0.3\linewidth}}{
        \cellcolor[rgb]{0.9,0.9,0.9}{
            \makecell[{{p{\linewidth}}}]{
                \texttt{\tiny{[GM$|$GM]}}
                \texttt{A moved to (1, 6)} \\
            }
        }
    }
    & & \\ \\

    \theutterance \stepcounter{utterance}  
    & & \multicolumn{4}{p{0.6\linewidth}}{
        \cellcolor[rgb]{0.9,0.9,0.9}{
            \makecell[{{p{\linewidth}}}]{
                \texttt{\tiny{[GM$\rangle$P2]}}
                \texttt{Vaše zpráva byla předána druhému hráči.} \\
\\ 
\\ 
\texttt{Právě hrajete kolo 6 z maximálně 28.} \\
\texttt{Dohromady jste nasbírali 7 z 16 pokut.} \\
\texttt{Druhý hráč přesunul objekt na své mřížce.} \\
\\ 
\\ 
\texttt{Jaký je váš další příkaz?} \\
            }
        }
    }
    & \\ \\

    \theutterance \stepcounter{utterance}  
    & & & \multicolumn{4}{p{0.6\linewidth}}{
        \cellcolor[rgb]{0.9,0.9,0.9}{
            \makecell[{{p{\linewidth}}}]{
                \texttt{\tiny{[GM$\langle$P2]}}
                \texttt{POHYB: A, (1, 6)} \\
            }
        }
    }
     \\ \\

    \theutterance \stepcounter{utterance}  
    & & & \multicolumn{2}{p{0.3\linewidth}}{
        \cellcolor[rgb]{0.9,0.9,0.9}{
            \makecell[{{p{\linewidth}}}]{
                \texttt{\tiny{[GM$|$GM]}}
                \texttt{POHYB: A, (1, 6)} \\
            }
        }
    }
    & & \\ \\

    \theutterance \stepcounter{utterance}  
    & & & \multicolumn{2}{p{0.3\linewidth}}{
        \cellcolor[rgb]{0.9,0.9,0.9}{
            \makecell[{{p{\linewidth}}}]{
                \texttt{\tiny{[GM$|$GM]}}
                \texttt{A moved to (1, 6)} \\
            }
        }
    }
    & & \\ \\

    \theutterance \stepcounter{utterance}  
    & & \multicolumn{4}{p{0.6\linewidth}}{
        \cellcolor[rgb]{0.9,0.9,0.9}{
            \makecell[{{p{\linewidth}}}]{
                \texttt{\tiny{[P1$\langle$GM]}}
                \texttt{Úspěšně přesunuto 'A' do (1,6). Vaše aktualizovaná mřížka vypadá takto:} \\
\texttt{```} \\
\texttt{ 1234567} \\
\texttt{╔═╤═══╤═╗} \\
\texttt{║◌│◌DT│◌║ 1} \\
\texttt{╟─┼──┬┘◌║ 2} \\
\texttt{║◌│◌◌│◌◌║ 3} \\
\texttt{╟─┴──┘◌O║ 4} \\
\texttt{║◌M◌P◌◌◌║ 5} \\
\texttt{║A◌◌◌◌◌◌║ 6} \\
\texttt{║S◌◌◌◌◌◌║ 7} \\
\texttt{╚═══════╝} \\
\texttt{```} \\
\texttt{Právě hrajete kolo 6 z maximálně 28.} \\
\texttt{Dohromady jste nasbírali 7 z 16 pokut.} \\
\texttt{Druhý hráč přesunul objekt na své mřížce.} \\
\\ 
\\ 
\texttt{Jaký je váš další příkaz?} \\
            }
        }
    }
    & \\ \\

    \theutterance \stepcounter{utterance}  
    & \multicolumn{4}{p{0.6\linewidth}}{
        \cellcolor[rgb]{0.9,0.9,0.9}{
            \makecell[{{p{\linewidth}}}]{
                \texttt{\tiny{[P1$\rangle$GM]}}
                \texttt{POHYB: M, (2, 6)} \\
            }
        }
    }
    & & \\ \\

    \theutterance \stepcounter{utterance}  
    & & & \multicolumn{2}{p{0.3\linewidth}}{
        \cellcolor[rgb]{0.9,0.9,0.9}{
            \makecell[{{p{\linewidth}}}]{
                \texttt{\tiny{[GM$|$GM]}}
                \texttt{POHYB: M, (2, 6)} \\
            }
        }
    }
    & & \\ \\

    \theutterance \stepcounter{utterance}  
    & & & \multicolumn{2}{p{0.3\linewidth}}{
        \cellcolor[rgb]{0.9,0.9,0.9}{
            \makecell[{{p{\linewidth}}}]{
                \texttt{\tiny{[GM$|$GM]}}
                \texttt{M moved to (2, 6)} \\
            }
        }
    }
    & & \\ \\

    \theutterance \stepcounter{utterance}  
    & & \multicolumn{4}{p{0.6\linewidth}}{
        \cellcolor[rgb]{0.9,0.9,0.9}{
            \makecell[{{p{\linewidth}}}]{
                \texttt{\tiny{[GM$\rangle$P2]}}
                \texttt{Úspěšně přesunuto 'A' do (1,6). Vaše aktualizovaná mřížka vypadá takto:} \\
\texttt{```} \\
\texttt{ 1234567} \\
\texttt{╔═╤═══╤═╗} \\
\texttt{║◌│◌◌◌│◌║ 1} \\
\texttt{╟─┼──┬┘◌║ 2} \\
\texttt{║◌│◌◌│◌◌║ 3} \\
\texttt{╟─┴──┘◌S║ 4} \\
\texttt{║◌◌M◌T◌◌║ 5} \\
\texttt{║AP◌◌◌◌◌║ 6} \\
\texttt{║◌◌O◌D◌◌║ 7} \\
\texttt{╚═══════╝} \\
\texttt{```} \\
\texttt{Právě hrajete kolo 7 z maximálně 28.} \\
\texttt{Dohromady jste nasbírali 7 z 16 pokut.} \\
\texttt{Druhý hráč přesunul objekt na své mřížce.} \\
\\ 
\\ 
\texttt{Jaký je váš další příkaz?} \\
            }
        }
    }
    & \\ \\

    \theutterance \stepcounter{utterance}  
    & & & \multicolumn{4}{p{0.6\linewidth}}{
        \cellcolor[rgb]{0.9,0.9,0.9}{
            \makecell[{{p{\linewidth}}}]{
                \texttt{\tiny{[GM$\langle$P2]}}
                \texttt{POHYB: M, (2, 6)} \\
            }
        }
    }
     \\ \\

    \theutterance \stepcounter{utterance}  
    & & & \multicolumn{2}{p{0.3\linewidth}}{
        \cellcolor[rgb]{0.9,0.9,0.9}{
            \makecell[{{p{\linewidth}}}]{
                \texttt{\tiny{[GM$|$GM]}}
                \texttt{POHYB: M, (2, 6)} \\
            }
        }
    }
    & & \\ \\

    \theutterance \stepcounter{utterance}  
    & & & \multicolumn{2}{p{0.3\linewidth}}{
        \cellcolor[rgb]{0.9,0.9,0.9}{
            \makecell[{{p{\linewidth}}}]{
                \texttt{\tiny{[GM$|$GM]}}
                \texttt{Penalizace: (2,6) není prázdný, ale obsahuje 'P'.} \\
\\ 
\texttt{Dohromady jste nasbírali 8 z 16 pokut.} \\
\texttt{Zkuste to prosím znovu!} \\
            }
        }
    }
    & & \\ \\

    \theutterance \stepcounter{utterance}  
    & & \multicolumn{4}{p{0.6\linewidth}}{
        \cellcolor[rgb]{0.9,0.9,0.9}{
            \makecell[{{p{\linewidth}}}]{
                \texttt{\tiny{[GM$\rangle$P2]}}
                \texttt{Penalizace: (2,6) není prázdný, ale obsahuje 'P'.} \\
\\ 
\texttt{Dohromady jste nasbírali 8 z 16 pokut.} \\
\texttt{Zkuste to prosím znovu!} \\
            }
        }
    }
    & \\ \\

    \theutterance \stepcounter{utterance}  
    & & & \multicolumn{4}{p{0.6\linewidth}}{
        \cellcolor[rgb]{0.9,0.9,0.9}{
            \makecell[{{p{\linewidth}}}]{
                \texttt{\tiny{[GM$\langle$P2]}}
                \texttt{POHYB: M, (3, 6)} \\
            }
        }
    }
     \\ \\

    \theutterance \stepcounter{utterance}  
    & & & \multicolumn{2}{p{0.3\linewidth}}{
        \cellcolor[rgb]{0.9,0.9,0.9}{
            \makecell[{{p{\linewidth}}}]{
                \texttt{\tiny{[GM$|$GM]}}
                \texttt{POHYB: M, (3, 6)} \\
            }
        }
    }
    & & \\ \\

    \theutterance \stepcounter{utterance}  
    & & & \multicolumn{2}{p{0.3\linewidth}}{
        \cellcolor[rgb]{0.9,0.9,0.9}{
            \makecell[{{p{\linewidth}}}]{
                \texttt{\tiny{[GM$|$GM]}}
                \texttt{M moved to (3, 6)} \\
            }
        }
    }
    & & \\ \\

    \theutterance \stepcounter{utterance}  
    & & \multicolumn{4}{p{0.6\linewidth}}{
        \cellcolor[rgb]{0.9,0.9,0.9}{
            \makecell[{{p{\linewidth}}}]{
                \texttt{\tiny{[P1$\langle$GM]}}
                \texttt{Úspěšně přesunuto 'M' do (2,6). Vaše aktualizovaná mřížka vypadá takto:} \\
\texttt{```} \\
\texttt{ 1234567} \\
\texttt{╔═╤═══╤═╗} \\
\texttt{║◌│◌DT│◌║ 1} \\
\texttt{╟─┼──┬┘◌║ 2} \\
\texttt{║◌│◌◌│◌◌║ 3} \\
\texttt{╟─┴──┘◌O║ 4} \\
\texttt{║◌◌◌P◌◌◌║ 5} \\
\texttt{║AM◌◌◌◌◌║ 6} \\
\texttt{║S◌◌◌◌◌◌║ 7} \\
\texttt{╚═══════╝} \\
\texttt{```} \\
\texttt{Právě hrajete kolo 7 z maximálně 28.} \\
\texttt{Dohromady jste nasbírali 8 z 16 pokut.} \\
\texttt{Druhý hráč přesunul objekt na své mřížce.} \\
\\ 
\\ 
\texttt{Jaký je váš další příkaz?} \\
            }
        }
    }
    & \\ \\

    \theutterance \stepcounter{utterance}  
    & \multicolumn{4}{p{0.6\linewidth}}{
        \cellcolor[rgb]{0.9,0.9,0.9}{
            \makecell[{{p{\linewidth}}}]{
                \texttt{\tiny{[P1$\rangle$GM]}}
                \texttt{POHYB: T, (3, 6)} \\
            }
        }
    }
    & & \\ \\

    \theutterance \stepcounter{utterance}  
    & & & \multicolumn{2}{p{0.3\linewidth}}{
        \cellcolor[rgb]{0.9,0.9,0.9}{
            \makecell[{{p{\linewidth}}}]{
                \texttt{\tiny{[GM$|$GM]}}
                \texttt{POHYB: T, (3, 6)} \\
            }
        }
    }
    & & \\ \\

    \theutterance \stepcounter{utterance}  
    & & & \multicolumn{2}{p{0.3\linewidth}}{
        \cellcolor[rgb]{0.9,0.9,0.9}{
            \makecell[{{p{\linewidth}}}]{
                \texttt{\tiny{[GM$|$GM]}}
                \texttt{T moved to (3, 6)} \\
            }
        }
    }
    & & \\ \\

    \theutterance \stepcounter{utterance}  
    & & \multicolumn{4}{p{0.6\linewidth}}{
        \cellcolor[rgb]{0.9,0.9,0.9}{
            \makecell[{{p{\linewidth}}}]{
                \texttt{\tiny{[GM$\rangle$P2]}}
                \texttt{Úspěšně přesunuto 'M' do (3,6). Vaše aktualizovaná mřížka vypadá takto:} \\
\texttt{```} \\
\texttt{ 1234567} \\
\texttt{╔═╤═══╤═╗} \\
\texttt{║◌│◌◌◌│◌║ 1} \\
\texttt{╟─┼──┬┘◌║ 2} \\
\texttt{║◌│◌◌│◌◌║ 3} \\
\texttt{╟─┴──┘◌S║ 4} \\
\texttt{║◌◌◌◌T◌◌║ 5} \\
\texttt{║APM◌◌◌◌║ 6} \\
\texttt{║◌◌O◌D◌◌║ 7} \\
\texttt{╚═══════╝} \\
\texttt{```} \\
\texttt{Právě hrajete kolo 8 z maximálně 28.} \\
\texttt{Dohromady jste nasbírali 8 z 16 pokut.} \\
\texttt{Druhý hráč přesunul objekt na své mřížce.} \\
\\ 
\\ 
\texttt{Jaký je váš další příkaz?} \\
            }
        }
    }
    & \\ \\

    \theutterance \stepcounter{utterance}  
    & & & \multicolumn{4}{p{0.6\linewidth}}{
        \cellcolor[rgb]{0.9,0.9,0.9}{
            \makecell[{{p{\linewidth}}}]{
                \texttt{\tiny{[GM$\langle$P2]}}
                \texttt{POHYB: T, (3, 6)} \\
            }
        }
    }
     \\ \\

    \theutterance \stepcounter{utterance}  
    & & & \multicolumn{2}{p{0.3\linewidth}}{
        \cellcolor[rgb]{0.9,0.9,0.9}{
            \makecell[{{p{\linewidth}}}]{
                \texttt{\tiny{[GM$|$GM]}}
                \texttt{POHYB: T, (3, 6)} \\
            }
        }
    }
    & & \\ \\

    \theutterance \stepcounter{utterance}  
    & & & \multicolumn{2}{p{0.3\linewidth}}{
        \cellcolor[rgb]{0.9,0.9,0.9}{
            \makecell[{{p{\linewidth}}}]{
                \texttt{\tiny{[GM$|$GM]}}
                \texttt{Penalizace: (3,6) není prázdný, ale obsahuje 'M'.} \\
\\ 
\texttt{Dohromady jste nasbírali 9 z 16 pokut.} \\
\texttt{Zkuste to prosím znovu!} \\
            }
        }
    }
    & & \\ \\

    \theutterance \stepcounter{utterance}  
    & & \multicolumn{4}{p{0.6\linewidth}}{
        \cellcolor[rgb]{0.9,0.9,0.9}{
            \makecell[{{p{\linewidth}}}]{
                \texttt{\tiny{[GM$\rangle$P2]}}
                \texttt{Penalizace: (3,6) není prázdný, ale obsahuje 'M'.} \\
\\ 
\texttt{Dohromady jste nasbírali 9 z 16 pokut.} \\
\texttt{Zkuste to prosím znovu!} \\
            }
        }
    }
    & \\ \\

    \theutterance \stepcounter{utterance}  
    & & & \multicolumn{4}{p{0.6\linewidth}}{
        \cellcolor[rgb]{0.9,0.9,0.9}{
            \makecell[{{p{\linewidth}}}]{
                \texttt{\tiny{[GM$\langle$P2]}}
                \texttt{POHYB: T, (4, 6)} \\
            }
        }
    }
     \\ \\

    \theutterance \stepcounter{utterance}  
    & & & \multicolumn{2}{p{0.3\linewidth}}{
        \cellcolor[rgb]{0.9,0.9,0.9}{
            \makecell[{{p{\linewidth}}}]{
                \texttt{\tiny{[GM$|$GM]}}
                \texttt{POHYB: T, (4, 6)} \\
            }
        }
    }
    & & \\ \\

    \theutterance \stepcounter{utterance}  
    & & & \multicolumn{2}{p{0.3\linewidth}}{
        \cellcolor[rgb]{0.9,0.9,0.9}{
            \makecell[{{p{\linewidth}}}]{
                \texttt{\tiny{[GM$|$GM]}}
                \texttt{T moved to (4, 6)} \\
            }
        }
    }
    & & \\ \\

    \theutterance \stepcounter{utterance}  
    & & \multicolumn{4}{p{0.6\linewidth}}{
        \cellcolor[rgb]{0.9,0.9,0.9}{
            \makecell[{{p{\linewidth}}}]{
                \texttt{\tiny{[P1$\langle$GM]}}
                \texttt{Úspěšně přesunuto 'T' do (3,6). Vaše aktualizovaná mřížka vypadá takto:} \\
\texttt{```} \\
\texttt{ 1234567} \\
\texttt{╔═╤═══╤═╗} \\
\texttt{║◌│◌D◌│◌║ 1} \\
\texttt{╟─┼──┬┘◌║ 2} \\
\texttt{║◌│◌◌│◌◌║ 3} \\
\texttt{╟─┴──┘◌O║ 4} \\
\texttt{║◌◌◌P◌◌◌║ 5} \\
\texttt{║AMT◌◌◌◌║ 6} \\
\texttt{║S◌◌◌◌◌◌║ 7} \\
\texttt{╚═══════╝} \\
\texttt{```} \\
\texttt{Právě hrajete kolo 8 z maximálně 28.} \\
\texttt{Dohromady jste nasbírali 9 z 16 pokut.} \\
\texttt{Druhý hráč přesunul objekt na své mřížce.} \\
\\ 
\\ 
\texttt{Jaký je váš další příkaz?} \\
            }
        }
    }
    & \\ \\

    \theutterance \stepcounter{utterance}  
    & \multicolumn{4}{p{0.6\linewidth}}{
        \cellcolor[rgb]{0.9,0.9,0.9}{
            \makecell[{{p{\linewidth}}}]{
                \texttt{\tiny{[P1$\rangle$GM]}}
                \texttt{POHYB: S, (4, 6)} \\
            }
        }
    }
    & & \\ \\

    \theutterance \stepcounter{utterance}  
    & & & \multicolumn{2}{p{0.3\linewidth}}{
        \cellcolor[rgb]{0.9,0.9,0.9}{
            \makecell[{{p{\linewidth}}}]{
                \texttt{\tiny{[GM$|$GM]}}
                \texttt{POHYB: S, (4, 6)} \\
            }
        }
    }
    & & \\ \\

    \theutterance \stepcounter{utterance}  
    & & & \multicolumn{2}{p{0.3\linewidth}}{
        \cellcolor[rgb]{0.9,0.9,0.9}{
            \makecell[{{p{\linewidth}}}]{
                \texttt{\tiny{[GM$|$GM]}}
                \texttt{S moved to (4, 6)} \\
            }
        }
    }
    & & \\ \\

    \theutterance \stepcounter{utterance}  
    & & \multicolumn{4}{p{0.6\linewidth}}{
        \cellcolor[rgb]{0.9,0.9,0.9}{
            \makecell[{{p{\linewidth}}}]{
                \texttt{\tiny{[GM$\rangle$P2]}}
                \texttt{Úspěšně přesunuto 'T' do (4,6). Vaše aktualizovaná mřížka vypadá takto:} \\
\texttt{```} \\
\texttt{ 1234567} \\
\texttt{╔═╤═══╤═╗} \\
\texttt{║◌│◌◌◌│◌║ 1} \\
\texttt{╟─┼──┬┘◌║ 2} \\
\texttt{║◌│◌◌│◌◌║ 3} \\
\texttt{╟─┴──┘◌S║ 4} \\
\texttt{║◌◌◌◌◌◌◌║ 5} \\
\texttt{║APMT◌◌◌║ 6} \\
\texttt{║◌◌O◌D◌◌║ 7} \\
\texttt{╚═══════╝} \\
\texttt{```} \\
\texttt{Právě hrajete kolo 9 z maximálně 28.} \\
\texttt{Dohromady jste nasbírali 9 z 16 pokut.} \\
\texttt{Druhý hráč přesunul objekt na své mřížce.} \\
\\ 
\\ 
\texttt{Jaký je váš další příkaz?} \\
            }
        }
    }
    & \\ \\

    \theutterance \stepcounter{utterance}  
    & & & \multicolumn{4}{p{0.6\linewidth}}{
        \cellcolor[rgb]{0.9,0.9,0.9}{
            \makecell[{{p{\linewidth}}}]{
                \texttt{\tiny{[GM$\langle$P2]}}
                \texttt{POHYB: P, (5, 6)} \\
            }
        }
    }
     \\ \\

    \theutterance \stepcounter{utterance}  
    & & & \multicolumn{2}{p{0.3\linewidth}}{
        \cellcolor[rgb]{0.9,0.9,0.9}{
            \makecell[{{p{\linewidth}}}]{
                \texttt{\tiny{[GM$|$GM]}}
                \texttt{POHYB: P, (5, 6)} \\
            }
        }
    }
    & & \\ \\

    \theutterance \stepcounter{utterance}  
    & & & \multicolumn{2}{p{0.3\linewidth}}{
        \cellcolor[rgb]{0.9,0.9,0.9}{
            \makecell[{{p{\linewidth}}}]{
                \texttt{\tiny{[GM$|$GM]}}
                \texttt{P moved to (5, 6)} \\
            }
        }
    }
    & & \\ \\

    \theutterance \stepcounter{utterance}  
    & & \multicolumn{4}{p{0.6\linewidth}}{
        \cellcolor[rgb]{0.9,0.9,0.9}{
            \makecell[{{p{\linewidth}}}]{
                \texttt{\tiny{[P1$\langle$GM]}}
                \texttt{Úspěšně přesunuto 'S' do (4,6). Vaše aktualizovaná mřížka vypadá takto:} \\
\texttt{```} \\
\texttt{ 1234567} \\
\texttt{╔═╤═══╤═╗} \\
\texttt{║◌│◌D◌│◌║ 1} \\
\texttt{╟─┼──┬┘◌║ 2} \\
\texttt{║◌│◌◌│◌◌║ 3} \\
\texttt{╟─┴──┘◌O║ 4} \\
\texttt{║◌◌◌P◌◌◌║ 5} \\
\texttt{║AMTS◌◌◌║ 6} \\
\texttt{║◌◌◌◌◌◌◌║ 7} \\
\texttt{╚═══════╝} \\
\texttt{```} \\
\texttt{Právě hrajete kolo 9 z maximálně 28.} \\
\texttt{Dohromady jste nasbírali 9 z 16 pokut.} \\
\texttt{Druhý hráč přesunul objekt na své mřížce.} \\
\\ 
\\ 
\texttt{Jaký je váš další příkaz?} \\
            }
        }
    }
    & \\ \\

    \theutterance \stepcounter{utterance}  
    & \multicolumn{4}{p{0.6\linewidth}}{
        \cellcolor[rgb]{0.9,0.9,0.9}{
            \makecell[{{p{\linewidth}}}]{
                \texttt{\tiny{[P1$\rangle$GM]}}
                \texttt{POHYB: P, (5, 6)} \\
            }
        }
    }
    & & \\ \\

    \theutterance \stepcounter{utterance}  
    & & & \multicolumn{2}{p{0.3\linewidth}}{
        \cellcolor[rgb]{0.9,0.9,0.9}{
            \makecell[{{p{\linewidth}}}]{
                \texttt{\tiny{[GM$|$GM]}}
                \texttt{POHYB: P, (5, 6)} \\
            }
        }
    }
    & & \\ \\

    \theutterance \stepcounter{utterance}  
    & & & \multicolumn{2}{p{0.3\linewidth}}{
        \cellcolor[rgb]{0.9,0.9,0.9}{
            \makecell[{{p{\linewidth}}}]{
                \texttt{\tiny{[GM$|$GM]}}
                \texttt{P moved to (5, 6)} \\
            }
        }
    }
    & & \\ \\

    \theutterance \stepcounter{utterance}  
    & & \multicolumn{4}{p{0.6\linewidth}}{
        \cellcolor[rgb]{0.9,0.9,0.9}{
            \makecell[{{p{\linewidth}}}]{
                \texttt{\tiny{[GM$\rangle$P2]}}
                \texttt{Úspěšně přesunuto 'P' do (5,6). Vaše aktualizovaná mřížka vypadá takto:} \\
\texttt{```} \\
\texttt{ 1234567} \\
\texttt{╔═╤═══╤═╗} \\
\texttt{║◌│◌◌◌│◌║ 1} \\
\texttt{╟─┼──┬┘◌║ 2} \\
\texttt{║◌│◌◌│◌◌║ 3} \\
\texttt{╟─┴──┘◌S║ 4} \\
\texttt{║◌◌◌◌◌◌◌║ 5} \\
\texttt{║A◌MTP◌◌║ 6} \\
\texttt{║◌◌O◌D◌◌║ 7} \\
\texttt{╚═══════╝} \\
\texttt{```} \\
\texttt{Právě hrajete kolo 10 z maximálně 28.} \\
\texttt{Dohromady jste nasbírali 9 z 16 pokut.} \\
\texttt{Druhý hráč přesunul objekt na své mřížce.} \\
\\ 
\\ 
\texttt{Jaký je váš další příkaz?} \\
            }
        }
    }
    & \\ \\

    \theutterance \stepcounter{utterance}  
    & & & \multicolumn{4}{p{0.6\linewidth}}{
        \cellcolor[rgb]{0.9,0.9,0.9}{
            \makecell[{{p{\linewidth}}}]{
                \texttt{\tiny{[GM$\langle$P2]}}
                \texttt{POHYB: S, (4, 6)} \\
\\ 
\texttt{Wait, I need to check if (4,6) is empty.} \\
\texttt{Looking at the grid:} \\
\texttt{Row 6: `║A◌MTP◌◌║ 6`} \\
\texttt{(1,6)=A, (2,6)=◌, (3,6)=M, (4,6)=T, (5,6)=P, (6,6)=◌, (7,6)=◌.} \\
\texttt{(4,6) contains 'T'. I cannot move S there.} \\
\\ 
\texttt{Let's re{-}evaluate the target for S. The target was S(4,6).} \\
\texttt{Current positions in row 6: A(1,6), M(3,6), T(4,6), P(5,6).} \\
\texttt{I need to move T out of (4,6) first.} \\
\\ 
\texttt{Let's move T to (6,6).} \\
\texttt{Wait, let me check if (6,6) is empty.} \\
\texttt{Row 6: `║A◌MTP◌◌║ 6`} \\
\texttt{(1,6)=A, (2,6)=◌, (3,6)=M, (4,6)=T, (5,6)=P, (6,6)=◌, (7,6)=◌.} \\
\texttt{Yes, (6,6) is empty.} \\
\\ 
\texttt{POHYB: T, (6, 6)} \\
            }
        }
    }
     \\ \\

    \theutterance \stepcounter{utterance}  
    & & & \multicolumn{2}{p{0.3\linewidth}}{
        \cellcolor[rgb]{0.9,0.9,0.9}{
            \makecell[{{p{\linewidth}}}]{
                \texttt{\tiny{[GM$|$GM]}}
                \texttt{POHYB: S, (4, 6)} \\
\\ 
\texttt{Wait, I need to check if (4,6) is empty.} \\
\texttt{Looking at the grid:} \\
\texttt{Row 6: `║A◌MTP◌◌║ 6`} \\
\texttt{(1,6)=A, (2,6)=◌, (3,6)=M, (4,6)=T, (5,6)=P, (6,6)=◌, (7,6)=◌.} \\
\texttt{(4,6) contains 'T'. I cannot move S there.} \\
\\ 
\texttt{Let's re{-}evaluate the target for S. The target was S(4,6).} \\
\texttt{Current positions in row 6: A(1,6), M(3,6), T(4,6), P(5,6).} \\
\texttt{I need to move T out of (4,6) first.} \\
\\ 
\texttt{Let's move T to (6,6).} \\
\texttt{Wait, let me check if (6,6) is empty.} \\
\texttt{Row 6: `║A◌MTP◌◌║ 6`} \\
\texttt{(1,6)=A, (2,6)=◌, (3,6)=M, (4,6)=T, (5,6)=P, (6,6)=◌, (7,6)=◌.} \\
\texttt{Yes, (6,6) is empty.} \\
\\ 
\texttt{POHYB: T, (6, 6)} \\
            }
        }
    }
    & & \\ \\

    \theutterance \stepcounter{utterance}  
    & & & \multicolumn{2}{p{0.3\linewidth}}{
        \cellcolor[rgb]{0.9,0.9,0.9}{
            \makecell[{{p{\linewidth}}}]{
                \texttt{\tiny{[GM$|$GM]}}
                \texttt{Invalid format: head is not empty: 'POHYB: S, (4, 6)} \\
\\ 
\texttt{Wait, I need to check if (4,6) is empty.} \\
\texttt{Looking at the grid:} \\
\texttt{Row 6: ║A◌MTP◌◌║ 6} \\
\texttt{(1,6)=A, (2,6)=◌, (3,6)=M, (4,6)=T, (5,6)=P, (6,6)=◌, (7,6)=◌.} \\
\texttt{(4,6) contains 'T'. I cannot move S there.} \\
\\ 
\texttt{Let's re{-}evaluate the target for S. The target was S(4,6).} \\
\texttt{Current positions in row 6: A(1,6), M(3,6), T(4,6), P(5,6).} \\
\texttt{I need to move T out of (4,6) first.} \\
\\ 
\texttt{Let's move T to (6,6).} \\
\texttt{Wait, let me check if (6,6) is empty.} \\
\texttt{Row 6: ║A◌MTP◌◌║ 6} \\
\texttt{(1,6)=A, (2,6)=◌, (3,6)=M, (4,6)=T, (5,6)=P, (6,6)=◌, (7,6)=◌.} \\
\texttt{Yes, (6,6) is empty.} \\
\\ 
\texttt{'} \\
            }
        }
    }
    & & \\ \\

    \theutterance \stepcounter{utterance}  
    & & \multicolumn{4}{p{0.6\linewidth}}{
        \cellcolor[rgb]{0.9,0.9,0.9}{
            \makecell[{{p{\linewidth}}}]{
                \texttt{\tiny{[GM$\rangle$P2]}}
                \texttt{Trest: Vaše zpráva nesmí obsahovat nic před příkazem!} \\
\\ 
\texttt{Ujistěte se, že vaše odpověď obsahuje pouze `ŘEKNI: <ZPRÁVA>` nebo `POHYB: <OBJEKT>, (<X>, <Y>)` a nic jiného!} \\
\\ 
\texttt{Dohromady jste nasbírali 10 z 16 pokut.} \\
\texttt{Zkuste to prosím znovu!} \\
            }
        }
    }
    & \\ \\

    \theutterance \stepcounter{utterance}  
    & & & \multicolumn{4}{p{0.6\linewidth}}{
        \cellcolor[rgb]{0.9,0.9,0.9}{
            \makecell[{{p{\linewidth}}}]{
                \texttt{\tiny{[GM$\langle$P2]}}
                \texttt{POHYB: T, (6, 6)} \\
            }
        }
    }
     \\ \\

    \theutterance \stepcounter{utterance}  
    & & & \multicolumn{2}{p{0.3\linewidth}}{
        \cellcolor[rgb]{0.9,0.9,0.9}{
            \makecell[{{p{\linewidth}}}]{
                \texttt{\tiny{[GM$|$GM]}}
                \texttt{POHYB: T, (6, 6)} \\
            }
        }
    }
    & & \\ \\

    \theutterance \stepcounter{utterance}  
    & & & \multicolumn{2}{p{0.3\linewidth}}{
        \cellcolor[rgb]{0.9,0.9,0.9}{
            \makecell[{{p{\linewidth}}}]{
                \texttt{\tiny{[GM$|$GM]}}
                \texttt{T moved to (6, 6)} \\
            }
        }
    }
    & & \\ \\

    \theutterance \stepcounter{utterance}  
    & & \multicolumn{4}{p{0.6\linewidth}}{
        \cellcolor[rgb]{0.9,0.9,0.9}{
            \makecell[{{p{\linewidth}}}]{
                \texttt{\tiny{[P1$\langle$GM]}}
                \texttt{Úspěšně přesunuto 'P' do (5,6). Vaše aktualizovaná mřížka vypadá takto:} \\
\texttt{```} \\
\texttt{ 1234567} \\
\texttt{╔═╤═══╤═╗} \\
\texttt{║◌│◌D◌│◌║ 1} \\
\texttt{╟─┼──┬┘◌║ 2} \\
\texttt{║◌│◌◌│◌◌║ 3} \\
\texttt{╟─┴──┘◌O║ 4} \\
\texttt{║◌◌◌◌◌◌◌║ 5} \\
\texttt{║AMTSP◌◌║ 6} \\
\texttt{║◌◌◌◌◌◌◌║ 7} \\
\texttt{╚═══════╝} \\
\texttt{```} \\
\texttt{Právě hrajete kolo 10 z maximálně 28.} \\
\texttt{Dohromady jste nasbírali 10 z 16 pokut.} \\
\texttt{Druhý hráč přesunul objekt na své mřížce.} \\
\\ 
\\ 
\texttt{Jaký je váš další příkaz?} \\
            }
        }
    }
    & \\ \\

    \theutterance \stepcounter{utterance}  
    & \multicolumn{4}{p{0.6\linewidth}}{
        \cellcolor[rgb]{0.9,0.9,0.9}{
            \makecell[{{p{\linewidth}}}]{
                \texttt{\tiny{[P1$\rangle$GM]}}
                \texttt{POHYB: O, (6, 6)} \\
            }
        }
    }
    & & \\ \\

    \theutterance \stepcounter{utterance}  
    & & & \multicolumn{2}{p{0.3\linewidth}}{
        \cellcolor[rgb]{0.9,0.9,0.9}{
            \makecell[{{p{\linewidth}}}]{
                \texttt{\tiny{[GM$|$GM]}}
                \texttt{POHYB: O, (6, 6)} \\
            }
        }
    }
    & & \\ \\

    \theutterance \stepcounter{utterance}  
    & & & \multicolumn{2}{p{0.3\linewidth}}{
        \cellcolor[rgb]{0.9,0.9,0.9}{
            \makecell[{{p{\linewidth}}}]{
                \texttt{\tiny{[GM$|$GM]}}
                \texttt{O moved to (6, 6)} \\
            }
        }
    }
    & & \\ \\

    \theutterance \stepcounter{utterance}  
    & & \multicolumn{4}{p{0.6\linewidth}}{
        \cellcolor[rgb]{0.9,0.9,0.9}{
            \makecell[{{p{\linewidth}}}]{
                \texttt{\tiny{[GM$\rangle$P2]}}
                \texttt{Úspěšně přesunuto 'T' do (6,6). Vaše aktualizovaná mřížka vypadá takto:} \\
\texttt{```} \\
\texttt{ 1234567} \\
\texttt{╔═╤═══╤═╗} \\
\texttt{║◌│◌◌◌│◌║ 1} \\
\texttt{╟─┼──┬┘◌║ 2} \\
\texttt{║◌│◌◌│◌◌║ 3} \\
\texttt{╟─┴──┘◌S║ 4} \\
\texttt{║◌◌◌◌◌◌◌║ 5} \\
\texttt{║A◌M◌PT◌║ 6} \\
\texttt{║◌◌O◌D◌◌║ 7} \\
\texttt{╚═══════╝} \\
\texttt{```} \\
\texttt{Právě hrajete kolo 11 z maximálně 28.} \\
\texttt{Dohromady jste nasbírali 10 z 16 pokut.} \\
\texttt{Druhý hráč přesunul objekt na své mřížce.} \\
\\ 
\\ 
\texttt{Jaký je váš další příkaz?} \\
            }
        }
    }
    & \\ \\

    \theutterance \stepcounter{utterance}  
    & & & \multicolumn{4}{p{0.6\linewidth}}{
        \cellcolor[rgb]{0.9,0.9,0.9}{
            \makecell[{{p{\linewidth}}}]{
                \texttt{\tiny{[GM$\langle$P2]}}
                \texttt{POHYB: S, (4, 6)} \\
            }
        }
    }
     \\ \\

    \theutterance \stepcounter{utterance}  
    & & & \multicolumn{2}{p{0.3\linewidth}}{
        \cellcolor[rgb]{0.9,0.9,0.9}{
            \makecell[{{p{\linewidth}}}]{
                \texttt{\tiny{[GM$|$GM]}}
                \texttt{POHYB: S, (4, 6)} \\
            }
        }
    }
    & & \\ \\

    \theutterance \stepcounter{utterance}  
    & & & \multicolumn{2}{p{0.3\linewidth}}{
        \cellcolor[rgb]{0.9,0.9,0.9}{
            \makecell[{{p{\linewidth}}}]{
                \texttt{\tiny{[GM$|$GM]}}
                \texttt{S moved to (4, 6)} \\
            }
        }
    }
    & & \\ \\

    \theutterance \stepcounter{utterance}  
    & & \multicolumn{4}{p{0.6\linewidth}}{
        \cellcolor[rgb]{0.9,0.9,0.9}{
            \makecell[{{p{\linewidth}}}]{
                \texttt{\tiny{[P1$\langle$GM]}}
                \texttt{Úspěšně přesunuto 'O' do (6,6). Vaše aktualizovaná mřížka vypadá takto:} \\
\texttt{```} \\
\texttt{ 1234567} \\
\texttt{╔═╤═══╤═╗} \\
\texttt{║◌│◌D◌│◌║ 1} \\
\texttt{╟─┼──┬┘◌║ 2} \\
\texttt{║◌│◌◌│◌◌║ 3} \\
\texttt{╟─┴──┘◌◌║ 4} \\
\texttt{║◌◌◌◌◌◌◌║ 5} \\
\texttt{║AMTSPO◌║ 6} \\
\texttt{║◌◌◌◌◌◌◌║ 7} \\
\texttt{╚═══════╝} \\
\texttt{```} \\
\texttt{Právě hrajete kolo 11 z maximálně 28.} \\
\texttt{Dohromady jste nasbírali 10 z 16 pokut.} \\
\texttt{Druhý hráč přesunul objekt na své mřížce.} \\
\\ 
\\ 
\texttt{Jaký je váš další příkaz?} \\
            }
        }
    }
    & \\ \\

    \theutterance \stepcounter{utterance}  
    & \multicolumn{4}{p{0.6\linewidth}}{
        \cellcolor[rgb]{0.9,0.9,0.9}{
            \makecell[{{p{\linewidth}}}]{
                \texttt{\tiny{[P1$\rangle$GM]}}
                \texttt{POHYB: D, (7, 6)} \\
            }
        }
    }
    & & \\ \\

    \theutterance \stepcounter{utterance}  
    & & & \multicolumn{2}{p{0.3\linewidth}}{
        \cellcolor[rgb]{0.9,0.9,0.9}{
            \makecell[{{p{\linewidth}}}]{
                \texttt{\tiny{[GM$|$GM]}}
                \texttt{POHYB: D, (7, 6)} \\
            }
        }
    }
    & & \\ \\

    \theutterance \stepcounter{utterance}  
    & & & \multicolumn{2}{p{0.3\linewidth}}{
        \cellcolor[rgb]{0.9,0.9,0.9}{
            \makecell[{{p{\linewidth}}}]{
                \texttt{\tiny{[GM$|$GM]}}
                \texttt{D moved to (7, 6)} \\
            }
        }
    }
    & & \\ \\

    \theutterance \stepcounter{utterance}  
    & & \multicolumn{4}{p{0.6\linewidth}}{
        \cellcolor[rgb]{0.9,0.9,0.9}{
            \makecell[{{p{\linewidth}}}]{
                \texttt{\tiny{[GM$\rangle$P2]}}
                \texttt{Úspěšně přesunuto 'S' do (4,6). Vaše aktualizovaná mřížka vypadá takto:} \\
\texttt{```} \\
\texttt{ 1234567} \\
\texttt{╔═╤═══╤═╗} \\
\texttt{║◌│◌◌◌│◌║ 1} \\
\texttt{╟─┼──┬┘◌║ 2} \\
\texttt{║◌│◌◌│◌◌║ 3} \\
\texttt{╟─┴──┘◌◌║ 4} \\
\texttt{║◌◌◌◌◌◌◌║ 5} \\
\texttt{║A◌MSPT◌║ 6} \\
\texttt{║◌◌O◌D◌◌║ 7} \\
\texttt{╚═══════╝} \\
\texttt{```} \\
\texttt{Právě hrajete kolo 12 z maximálně 28.} \\
\texttt{Dohromady jste nasbírali 10 z 16 pokut.} \\
\texttt{Druhý hráč přesunul objekt na své mřížce.} \\
\\ 
\\ 
\texttt{Jaký je váš další příkaz?} \\
            }
        }
    }
    & \\ \\

    \theutterance \stepcounter{utterance}  
    & & & \multicolumn{4}{p{0.6\linewidth}}{
        \cellcolor[rgb]{0.9,0.9,0.9}{
            \makecell[{{p{\linewidth}}}]{
                \texttt{\tiny{[GM$\langle$P2]}}
                \texttt{POHYB: M, (2, 6)} \\
            }
        }
    }
     \\ \\

    \theutterance \stepcounter{utterance}  
    & & & \multicolumn{2}{p{0.3\linewidth}}{
        \cellcolor[rgb]{0.9,0.9,0.9}{
            \makecell[{{p{\linewidth}}}]{
                \texttt{\tiny{[GM$|$GM]}}
                \texttt{POHYB: M, (2, 6)} \\
            }
        }
    }
    & & \\ \\

    \theutterance \stepcounter{utterance}  
    & & & \multicolumn{2}{p{0.3\linewidth}}{
        \cellcolor[rgb]{0.9,0.9,0.9}{
            \makecell[{{p{\linewidth}}}]{
                \texttt{\tiny{[GM$|$GM]}}
                \texttt{M moved to (2, 6)} \\
            }
        }
    }
    & & \\ \\

    \theutterance \stepcounter{utterance}  
    & & \multicolumn{4}{p{0.6\linewidth}}{
        \cellcolor[rgb]{0.9,0.9,0.9}{
            \makecell[{{p{\linewidth}}}]{
                \texttt{\tiny{[P1$\langle$GM]}}
                \texttt{Úspěšně přesunuto 'D' do (7,6). Vaše aktualizovaná mřížka vypadá takto:} \\
\texttt{```} \\
\texttt{ 1234567} \\
\texttt{╔═╤═══╤═╗} \\
\texttt{║◌│◌◌◌│◌║ 1} \\
\texttt{╟─┼──┬┘◌║ 2} \\
\texttt{║◌│◌◌│◌◌║ 3} \\
\texttt{╟─┴──┘◌◌║ 4} \\
\texttt{║◌◌◌◌◌◌◌║ 5} \\
\texttt{║AMTSPOD║ 6} \\
\texttt{║◌◌◌◌◌◌◌║ 7} \\
\texttt{╚═══════╝} \\
\texttt{```} \\
\texttt{Právě hrajete kolo 12 z maximálně 28.} \\
\texttt{Dohromady jste nasbírali 10 z 16 pokut.} \\
\texttt{Druhý hráč přesunul objekt na své mřížce.} \\
\\ 
\\ 
\texttt{Jaký je váš další příkaz?} \\
            }
        }
    }
    & \\ \\

    \theutterance \stepcounter{utterance}  
    & \multicolumn{4}{p{0.6\linewidth}}{
        \cellcolor[rgb]{0.9,0.9,0.9}{
            \makecell[{{p{\linewidth}}}]{
                \texttt{\tiny{[P1$\rangle$GM]}}
                \texttt{ŘEKNI: hotovo?} \\
            }
        }
    }
    & & \\ \\

    \theutterance \stepcounter{utterance}  
    & & & \multicolumn{2}{p{0.3\linewidth}}{
        \cellcolor[rgb]{0.9,0.9,0.9}{
            \makecell[{{p{\linewidth}}}]{
                \texttt{\tiny{[GM$|$GM]}}
                \texttt{ŘEKNI: hotovo?} \\
            }
        }
    }
    & & \\ \\

    \theutterance \stepcounter{utterance}  
    & & \multicolumn{4}{p{0.6\linewidth}}{
        \cellcolor[rgb]{0.9,0.9,0.9}{
            \makecell[{{p{\linewidth}}}]{
                \texttt{\tiny{[GM$\rangle$P2]}}
                \texttt{Úspěšně přesunuto 'M' do (2,6). Vaše aktualizovaná mřížka vypadá takto:} \\
\texttt{```} \\
\texttt{ 1234567} \\
\texttt{╔═╤═══╤═╗} \\
\texttt{║◌│◌◌◌│◌║ 1} \\
\texttt{╟─┼──┬┘◌║ 2} \\
\texttt{║◌│◌◌│◌◌║ 3} \\
\texttt{╟─┴──┘◌◌║ 4} \\
\texttt{║◌◌◌◌◌◌◌║ 5} \\
\texttt{║AM◌SPT◌║ 6} \\
\texttt{║◌◌O◌D◌◌║ 7} \\
\texttt{╚═══════╝} \\
\texttt{```} \\
\texttt{Právě hrajete kolo 13 z maximálně 28.} \\
\texttt{Dohromady jste nasbírali 10 z 16 pokut.} \\
\texttt{Druhý hráč poslal tuto zprávu:} \\
\\ 
\\ 
\texttt{"hotovo?"} \\
\\ 
\\ 
\texttt{Jaký je váš další příkaz?} \\
            }
        }
    }
    & \\ \\

    \theutterance \stepcounter{utterance}  
    & & & \multicolumn{4}{p{0.6\linewidth}}{
        \cellcolor[rgb]{0.9,0.9,0.9}{
            \makecell[{{p{\linewidth}}}]{
                \texttt{\tiny{[GM$\langle$P2]}}
                \texttt{POHYB: T, (3, 6)} \\
            }
        }
    }
     \\ \\

    \theutterance \stepcounter{utterance}  
    & & & \multicolumn{2}{p{0.3\linewidth}}{
        \cellcolor[rgb]{0.9,0.9,0.9}{
            \makecell[{{p{\linewidth}}}]{
                \texttt{\tiny{[GM$|$GM]}}
                \texttt{POHYB: T, (3, 6)} \\
            }
        }
    }
    & & \\ \\

    \theutterance \stepcounter{utterance}  
    & & & \multicolumn{2}{p{0.3\linewidth}}{
        \cellcolor[rgb]{0.9,0.9,0.9}{
            \makecell[{{p{\linewidth}}}]{
                \texttt{\tiny{[GM$|$GM]}}
                \texttt{T moved to (3, 6)} \\
            }
        }
    }
    & & \\ \\

    \theutterance \stepcounter{utterance}  
    & & \multicolumn{4}{p{0.6\linewidth}}{
        \cellcolor[rgb]{0.9,0.9,0.9}{
            \makecell[{{p{\linewidth}}}]{
                \texttt{\tiny{[P1$\langle$GM]}}
                \texttt{Vaše zpráva byla předána druhému hráči.} \\
\\ 
\\ 
\texttt{Právě hrajete kolo 13 z maximálně 28.} \\
\texttt{Dohromady jste nasbírali 10 z 16 pokut.} \\
\texttt{Druhý hráč přesunul objekt na své mřížce.} \\
\\ 
\\ 
\texttt{Jaký je váš další příkaz?} \\
            }
        }
    }
    & \\ \\

    \theutterance \stepcounter{utterance}  
    & \multicolumn{4}{p{0.6\linewidth}}{
        \cellcolor[rgb]{0.9,0.9,0.9}{
            \makecell[{{p{\linewidth}}}]{
                \texttt{\tiny{[P1$\rangle$GM]}}
                \texttt{ŘEKNI: hotovo!} \\
            }
        }
    }
    & & \\ \\

    \theutterance \stepcounter{utterance}  
    & & & \multicolumn{2}{p{0.3\linewidth}}{
        \cellcolor[rgb]{0.9,0.9,0.9}{
            \makecell[{{p{\linewidth}}}]{
                \texttt{\tiny{[GM$|$GM]}}
                \texttt{ŘEKNI: hotovo!} \\
            }
        }
    }
    & & \\ \\

    \theutterance \stepcounter{utterance}  
    & & \multicolumn{4}{p{0.6\linewidth}}{
        \cellcolor[rgb]{0.9,0.9,0.9}{
            \makecell[{{p{\linewidth}}}]{
                \texttt{\tiny{[GM$\rangle$P2]}}
                \texttt{Úspěšně přesunuto 'T' do (3,6). Vaše aktualizovaná mřížka vypadá takto:} \\
\texttt{```} \\
\texttt{ 1234567} \\
\texttt{╔═╤═══╤═╗} \\
\texttt{║◌│◌◌◌│◌║ 1} \\
\texttt{╟─┼──┬┘◌║ 2} \\
\texttt{║◌│◌◌│◌◌║ 3} \\
\texttt{╟─┴──┘◌◌║ 4} \\
\texttt{║◌◌◌◌◌◌◌║ 5} \\
\texttt{║AMTSP◌◌║ 6} \\
\texttt{║◌◌O◌D◌◌║ 7} \\
\texttt{╚═══════╝} \\
\texttt{```} \\
\texttt{Právě hrajete kolo 14 z maximálně 28.} \\
\texttt{Dohromady jste nasbírali 10 z 16 pokut.} \\
\texttt{Druhý hráč poslal tuto zprávu:} \\
\\ 
\\ 
\texttt{"hotovo!"} \\
\\ 
\\ 
\texttt{Jaký je váš další příkaz?} \\
            }
        }
    }
    & \\ \\

    \theutterance \stepcounter{utterance}  
    & & & \multicolumn{4}{p{0.6\linewidth}}{
        \cellcolor[rgb]{0.9,0.9,0.9}{
            \makecell[{{p{\linewidth}}}]{
                \texttt{\tiny{[GM$\langle$P2]}}
                \texttt{POHYB: P, (5, 6)} \\
            }
        }
    }
     \\ \\

    \theutterance \stepcounter{utterance}  
    & & & \multicolumn{2}{p{0.3\linewidth}}{
        \cellcolor[rgb]{0.9,0.9,0.9}{
            \makecell[{{p{\linewidth}}}]{
                \texttt{\tiny{[GM$|$GM]}}
                \texttt{POHYB: P, (5, 6)} \\
            }
        }
    }
    & & \\ \\

    \theutterance \stepcounter{utterance}  
    & & & \multicolumn{2}{p{0.3\linewidth}}{
        \cellcolor[rgb]{0.9,0.9,0.9}{
            \makecell[{{p{\linewidth}}}]{
                \texttt{\tiny{[GM$|$GM]}}
                \texttt{Penalizace: (5,6) není prázdný, ale obsahuje 'P'.} \\
\\ 
\texttt{Dohromady jste nasbírali 11 z 16 pokut.} \\
\texttt{Zkuste to prosím znovu!} \\
            }
        }
    }
    & & \\ \\

    \theutterance \stepcounter{utterance}  
    & & \multicolumn{4}{p{0.6\linewidth}}{
        \cellcolor[rgb]{0.9,0.9,0.9}{
            \makecell[{{p{\linewidth}}}]{
                \texttt{\tiny{[GM$\rangle$P2]}}
                \texttt{Penalizace: (5,6) není prázdný, ale obsahuje 'P'.} \\
\\ 
\texttt{Dohromady jste nasbírali 11 z 16 pokut.} \\
\texttt{Zkuste to prosím znovu!} \\
            }
        }
    }
    & \\ \\

    \theutterance \stepcounter{utterance}  
    & & & \multicolumn{4}{p{0.6\linewidth}}{
        \cellcolor[rgb]{0.9,0.9,0.9}{
            \makecell[{{p{\linewidth}}}]{
                \texttt{\tiny{[GM$\langle$P2]}}
                \texttt{POHYB: P, (5, 6)} \\
\\ 
\texttt{Wait, I see the error. Looking at the grid:} \\
\texttt{Row 6: `║AMTSP◌◌║ 6`} \\
\texttt{(1,6)=A, (2,6)=M, (3,6)=T, (4,6)=S, (5,6)=P, (6,6)=◌, (7,6)=◌.} \\
\texttt{The object 'P' is already at (5,6). My target for P was (5,6).} \\
\texttt{I should check the other objects.} \\
\texttt{Target: A(1,6), M(2,6), T(3,6), S(4,6), P(5,6), O(6,6), D(7,6).} \\
\\ 
\texttt{Current positions:} \\
\texttt{A is at (1,6) {-} OK} \\
\texttt{M is at (2,6) {-} OK} \\
\texttt{T is at (3,6) {-} OK} \\
\texttt{S is at (4,6) {-} OK} \\
\texttt{P is at (5,6) {-} OK} \\
\texttt{O is at (3,7) {-} Needs to move to (6,6)} \\
\texttt{D is at (5,7) {-} Needs to move to (7,6)} \\
\\ 
\texttt{I will move O to (6,6).} \\
\texttt{First, I must check if (6,6) is empty.} \\
\texttt{Row 6: `║AMTSP◌◌║ 6` {-}> (6,6) is `◌`. Yes.} \\
\\ 
\texttt{POHYB: O, (6, 6)} \\
            }
        }
    }
     \\ \\

    \theutterance \stepcounter{utterance}  
    & & & \multicolumn{2}{p{0.3\linewidth}}{
        \cellcolor[rgb]{0.9,0.9,0.9}{
            \makecell[{{p{\linewidth}}}]{
                \texttt{\tiny{[GM$|$GM]}}
                \texttt{POHYB: P, (5, 6)} \\
\\ 
\texttt{Wait, I see the error. Looking at the grid:} \\
\texttt{Row 6: `║AMTSP◌◌║ 6`} \\
\texttt{(1,6)=A, (2,6)=M, (3,6)=T, (4,6)=S, (5,6)=P, (6,6)=◌, (7,6)=◌.} \\
\texttt{The object 'P' is already at (5,6). My target for P was (5,6).} \\
\texttt{I should check the other objects.} \\
\texttt{Target: A(1,6), M(2,6), T(3,6), S(4,6), P(5,6), O(6,6), D(7,6).} \\
\\ 
\texttt{Current positions:} \\
\texttt{A is at (1,6) {-} OK} \\
\texttt{M is at (2,6) {-} OK} \\
\texttt{T is at (3,6) {-} OK} \\
\texttt{S is at (4,6) {-} OK} \\
\texttt{P is at (5,6) {-} OK} \\
\texttt{O is at (3,7) {-} Needs to move to (6,6)} \\
\texttt{D is at (5,7) {-} Needs to move to (7,6)} \\
\\ 
\texttt{I will move O to (6,6).} \\
\texttt{First, I must check if (6,6) is empty.} \\
\texttt{Row 6: `║AMTSP◌◌║ 6` {-}> (6,6) is `◌`. Yes.} \\
\\ 
\texttt{POHYB: O, (6, 6)} \\
            }
        }
    }
    & & \\ \\

    \theutterance \stepcounter{utterance}  
    & & & \multicolumn{2}{p{0.3\linewidth}}{
        \cellcolor[rgb]{0.9,0.9,0.9}{
            \makecell[{{p{\linewidth}}}]{
                \texttt{\tiny{[GM$|$GM]}}
                \texttt{Invalid format: head is not empty: 'POHYB: P, (5, 6)} \\
\\ 
\texttt{Wait, I see the error. Looking at the grid:} \\
\texttt{Row 6: ║AMTSP◌◌║ 6} \\
\texttt{(1,6)=A, (2,6)=M, (3,6)=T, (4,6)=S, (5,6)=P, (6,6)=◌, (7,6)=◌.} \\
\texttt{The object 'P' is already at (5,6). My target for P was (5,6).} \\
\texttt{I should check the other objects.} \\
\texttt{Target: A(1,6), M(2,6), T(3,6), S(4,6), P(5,6), O(6,6), D(7,6).} \\
\\ 
\texttt{Current positions:} \\
\texttt{A is at (1,6) {-} OK} \\
\texttt{M is at (2,6) {-} OK} \\
\texttt{T is at (3,6) {-} OK} \\
\texttt{S is at (4,6) {-} OK} \\
\texttt{P is at (5,6) {-} OK} \\
\texttt{O is at (3,7) {-} Needs to move to (6,6)} \\
\texttt{D is at (5,7) {-} Needs to move to (7,6)} \\
\\ 
\texttt{I will move O to (6,6).} \\
\texttt{First, I must check if (6,6) is empty.} \\
\texttt{Row 6: ║AMTSP◌◌║ 6 {-}> (6,6) is ◌. Yes.} \\
\\ 
\texttt{'} \\
            }
        }
    }
    & & \\ \\

    \theutterance \stepcounter{utterance}  
    & & \multicolumn{4}{p{0.6\linewidth}}{
        \cellcolor[rgb]{0.9,0.9,0.9}{
            \makecell[{{p{\linewidth}}}]{
                \texttt{\tiny{[GM$\rangle$P2]}}
                \texttt{Trest: Vaše zpráva nesmí obsahovat nic před příkazem!} \\
\\ 
\texttt{Ujistěte se, že vaše odpověď obsahuje pouze `ŘEKNI: <ZPRÁVA>` nebo `POHYB: <OBJEKT>, (<X>, <Y>)` a nic jiného!} \\
\\ 
\texttt{Dohromady jste nasbírali 12 z 16 pokut.} \\
\texttt{Zkuste to prosím znovu!} \\
            }
        }
    }
    & \\ \\

    \theutterance \stepcounter{utterance}  
    & & & \multicolumn{4}{p{0.6\linewidth}}{
        \cellcolor[rgb]{0.9,0.9,0.9}{
            \makecell[{{p{\linewidth}}}]{
                \texttt{\tiny{[GM$\langle$P2]}}
                \texttt{POHYB: O, (6, 6)} \\
            }
        }
    }
     \\ \\

    \theutterance \stepcounter{utterance}  
    & & & \multicolumn{2}{p{0.3\linewidth}}{
        \cellcolor[rgb]{0.9,0.9,0.9}{
            \makecell[{{p{\linewidth}}}]{
                \texttt{\tiny{[GM$|$GM]}}
                \texttt{POHYB: O, (6, 6)} \\
            }
        }
    }
    & & \\ \\

    \theutterance \stepcounter{utterance}  
    & & & \multicolumn{2}{p{0.3\linewidth}}{
        \cellcolor[rgb]{0.9,0.9,0.9}{
            \makecell[{{p{\linewidth}}}]{
                \texttt{\tiny{[GM$|$GM]}}
                \texttt{O moved to (6, 6)} \\
            }
        }
    }
    & & \\ \\

    \theutterance \stepcounter{utterance}  
    & & \multicolumn{4}{p{0.6\linewidth}}{
        \cellcolor[rgb]{0.9,0.9,0.9}{
            \makecell[{{p{\linewidth}}}]{
                \texttt{\tiny{[P1$\langle$GM]}}
                \texttt{Vaše zpráva byla předána druhému hráči.} \\
\\ 
\\ 
\texttt{Právě hrajete kolo 14 z maximálně 28.} \\
\texttt{Dohromady jste nasbírali 12 z 16 pokut.} \\
\texttt{Druhý hráč přesunul objekt na své mřížce.} \\
\\ 
\\ 
\texttt{Jaký je váš další příkaz?} \\
            }
        }
    }
    & \\ \\

    \theutterance \stepcounter{utterance}  
    & \multicolumn{4}{p{0.6\linewidth}}{
        \cellcolor[rgb]{0.9,0.9,0.9}{
            \makecell[{{p{\linewidth}}}]{
                \texttt{\tiny{[P1$\rangle$GM]}}
                \texttt{ŘEKNI: hotovo!} \\
            }
        }
    }
    & & \\ \\

    \theutterance \stepcounter{utterance}  
    & & & \multicolumn{2}{p{0.3\linewidth}}{
        \cellcolor[rgb]{0.9,0.9,0.9}{
            \makecell[{{p{\linewidth}}}]{
                \texttt{\tiny{[GM$|$GM]}}
                \texttt{ŘEKNI: hotovo!} \\
            }
        }
    }
    & & \\ \\

    \theutterance \stepcounter{utterance}  
    & & \multicolumn{4}{p{0.6\linewidth}}{
        \cellcolor[rgb]{0.9,0.9,0.9}{
            \makecell[{{p{\linewidth}}}]{
                \texttt{\tiny{[GM$\rangle$P2]}}
                \texttt{Úspěšně přesunuto 'O' do (6,6). Vaše aktualizovaná mřížka vypadá takto:} \\
\texttt{```} \\
\texttt{ 1234567} \\
\texttt{╔═╤═══╤═╗} \\
\texttt{║◌│◌◌◌│◌║ 1} \\
\texttt{╟─┼──┬┘◌║ 2} \\
\texttt{║◌│◌◌│◌◌║ 3} \\
\texttt{╟─┴──┘◌◌║ 4} \\
\texttt{║◌◌◌◌◌◌◌║ 5} \\
\texttt{║AMTSPO◌║ 6} \\
\texttt{║◌◌◌◌D◌◌║ 7} \\
\texttt{╚═══════╝} \\
\texttt{```} \\
\texttt{Právě hrajete kolo 15 z maximálně 28.} \\
\texttt{Dohromady jste nasbírali 12 z 16 pokut.} \\
\texttt{Druhý hráč poslal tuto zprávu:} \\
\\ 
\\ 
\texttt{"hotovo!"} \\
\\ 
\\ 
\texttt{Jaký je váš další příkaz?} \\
            }
        }
    }
    & \\ \\

    \theutterance \stepcounter{utterance}  
    & & & \multicolumn{4}{p{0.6\linewidth}}{
        \cellcolor[rgb]{0.9,0.9,0.9}{
            \makecell[{{p{\linewidth}}}]{
                \texttt{\tiny{[GM$\langle$P2]}}
                \texttt{POHYB: D, (7, 6)} \\
            }
        }
    }
     \\ \\

    \theutterance \stepcounter{utterance}  
    & & & \multicolumn{2}{p{0.3\linewidth}}{
        \cellcolor[rgb]{0.9,0.9,0.9}{
            \makecell[{{p{\linewidth}}}]{
                \texttt{\tiny{[GM$|$GM]}}
                \texttt{POHYB: D, (7, 6)} \\
            }
        }
    }
    & & \\ \\

    \theutterance \stepcounter{utterance}  
    & & & \multicolumn{2}{p{0.3\linewidth}}{
        \cellcolor[rgb]{0.9,0.9,0.9}{
            \makecell[{{p{\linewidth}}}]{
                \texttt{\tiny{[GM$|$GM]}}
                \texttt{D moved to (7, 6)} \\
            }
        }
    }
    & & \\ \\

    \theutterance \stepcounter{utterance}  
    & & \multicolumn{4}{p{0.6\linewidth}}{
        \cellcolor[rgb]{0.9,0.9,0.9}{
            \makecell[{{p{\linewidth}}}]{
                \texttt{\tiny{[P1$\langle$GM]}}
                \texttt{Vaše zpráva byla předána druhému hráči.} \\
\\ 
\\ 
\texttt{Právě hrajete kolo 15 z maximálně 28.} \\
\texttt{Dohromady jste nasbírali 12 z 16 pokut.} \\
\texttt{Druhý hráč přesunul objekt na své mřížce.} \\
\\ 
\\ 
\texttt{Jaký je váš další příkaz?} \\
            }
        }
    }
    & \\ \\

    \theutterance \stepcounter{utterance}  
    & \multicolumn{4}{p{0.6\linewidth}}{
        \cellcolor[rgb]{0.9,0.9,0.9}{
            \makecell[{{p{\linewidth}}}]{
                \texttt{\tiny{[P1$\rangle$GM]}}
                \texttt{ŘEKNI: hotovo!} \\
            }
        }
    }
    & & \\ \\

    \theutterance \stepcounter{utterance}  
    & & & \multicolumn{2}{p{0.3\linewidth}}{
        \cellcolor[rgb]{0.9,0.9,0.9}{
            \makecell[{{p{\linewidth}}}]{
                \texttt{\tiny{[GM$|$GM]}}
                \texttt{ŘEKNI: hotovo!} \\
            }
        }
    }
    & & \\ \\

    \theutterance \stepcounter{utterance}  
    & & \multicolumn{4}{p{0.6\linewidth}}{
        \cellcolor[rgb]{0.9,0.9,0.9}{
            \makecell[{{p{\linewidth}}}]{
                \texttt{\tiny{[GM$\rangle$P2]}}
                \texttt{Úspěšně přesunuto 'D' do (7,6). Vaše aktualizovaná mřížka vypadá takto:} \\
\texttt{```} \\
\texttt{ 1234567} \\
\texttt{╔═╤═══╤═╗} \\
\texttt{║◌│◌◌◌│◌║ 1} \\
\texttt{╟─┼──┬┘◌║ 2} \\
\texttt{║◌│◌◌│◌◌║ 3} \\
\texttt{╟─┴──┘◌◌║ 4} \\
\texttt{║◌◌◌◌◌◌◌║ 5} \\
\texttt{║AMTSPOD║ 6} \\
\texttt{║◌◌◌◌◌◌◌║ 7} \\
\texttt{╚═══════╝} \\
\texttt{```} \\
\texttt{Právě hrajete kolo 16 z maximálně 28.} \\
\texttt{Dohromady jste nasbírali 12 z 16 pokut.} \\
\texttt{Druhý hráč poslal tuto zprávu:} \\
\\ 
\\ 
\texttt{"hotovo!"} \\
\\ 
\\ 
\texttt{Jaký je váš další příkaz?} \\
            }
        }
    }
    & \\ \\

    \theutterance \stepcounter{utterance}  
    & & & \multicolumn{4}{p{0.6\linewidth}}{
        \cellcolor[rgb]{0.9,0.9,0.9}{
            \makecell[{{p{\linewidth}}}]{
                \texttt{\tiny{[GM$\langle$P2]}}
                \texttt{ŘEKNI: hotovo!} \\
            }
        }
    }
     \\ \\

    \theutterance \stepcounter{utterance}  
    & & & \multicolumn{2}{p{0.3\linewidth}}{
        \cellcolor[rgb]{0.9,0.9,0.9}{
            \makecell[{{p{\linewidth}}}]{
                \texttt{\tiny{[GM$|$GM]}}
                \texttt{ŘEKNI: hotovo!} \\
            }
        }
    }
    & & \\ \\

    \theutterance \stepcounter{utterance}  
    & & & \multicolumn{2}{p{0.3\linewidth}}{
        \cellcolor[rgb]{0.9,0.9,0.9}{
            \makecell[{{p{\linewidth}}}]{
                \texttt{\tiny{[GM$|$GM]}}
                \texttt{true} \\
            }
        }
    }
    & & \\ \\

    \theutterance \stepcounter{utterance}  
    & & & \multicolumn{2}{p{0.3\linewidth}}{
        \cellcolor[rgb]{0.9,0.9,0.9}{
            \makecell[{{p{\linewidth}}}]{
                \texttt{\tiny{[GM$|$GM]}}
                \texttt{Initial states:} \\
\texttt{```} \\
\texttt{Player 1:} \\
\texttt{ 1234567} \\
\texttt{╔═╤═══╤═╗} \\
\texttt{║◌│◌DT│◌║ 1} \\
\texttt{╟─┼──┬┘◌║ 2} \\
\texttt{║◌│◌◌│◌◌║ 3} \\
\texttt{╟─┴──┘◌O║ 4} \\
\texttt{║M◌◌P◌◌◌║ 5} \\
\texttt{║◌◌◌◌◌◌◌║ 6} \\
\texttt{║S◌◌◌◌◌A║ 7} \\
\texttt{╚═══════╝} \\
\\ 
\\ 
\texttt{Player 2:} \\
\texttt{ 1234567} \\
\texttt{╔═╤═══╤═╗} \\
\texttt{║◌│◌T◌│◌║ 1} \\
\texttt{╟─┼──┬┘◌║ 2} \\
\texttt{║A│M◌│◌◌║ 3} \\
\texttt{╟─┴──┘◌S║ 4} \\
\texttt{║◌◌◌◌◌◌◌║ 5} \\
\texttt{║◌P◌◌◌◌◌║ 6} \\
\texttt{║◌◌O◌D◌◌║ 7} \\
\texttt{╚═══════╝} \\
\\ 
\texttt{```} \\
            }
        }
    }
    & & \\ \\

    \theutterance \stepcounter{utterance}  
    & & & \multicolumn{2}{p{0.3\linewidth}}{
        \cellcolor[rgb]{0.9,0.9,0.9}{
            \makecell[{{p{\linewidth}}}]{
                \texttt{\tiny{[GM$|$GM]}}
                \texttt{End states:} \\
\texttt{```} \\
\texttt{Player 1:} \\
\texttt{ 1234567} \\
\texttt{╔═╤═══╤═╗} \\
\texttt{║◌│◌◌◌│◌║ 1} \\
\texttt{╟─┼──┬┘◌║ 2} \\
\texttt{║◌│◌◌│◌◌║ 3} \\
\texttt{╟─┴──┘◌◌║ 4} \\
\texttt{║◌◌◌◌◌◌◌║ 5} \\
\texttt{║AMTSPOD║ 6} \\
\texttt{║◌◌◌◌◌◌◌║ 7} \\
\texttt{╚═══════╝} \\
\\ 
\\ 
\texttt{Player 2:} \\
\texttt{ 1234567} \\
\texttt{╔═╤═══╤═╗} \\
\texttt{║◌│◌◌◌│◌║ 1} \\
\texttt{╟─┼──┬┘◌║ 2} \\
\texttt{║◌│◌◌│◌◌║ 3} \\
\texttt{╟─┴──┘◌◌║ 4} \\
\texttt{║◌◌◌◌◌◌◌║ 5} \\
\texttt{║AMTSPOD║ 6} \\
\texttt{║◌◌◌◌◌◌◌║ 7} \\
\texttt{╚═══════╝} \\
\\ 
\texttt{```} \\
            }
        }
    }
    & & \\ \\

    \theutterance \stepcounter{utterance}  
    & & & \multicolumn{2}{p{0.3\linewidth}}{
        \cellcolor[rgb]{0.9,0.9,0.9}{
            \makecell[{{p{\linewidth}}}]{
                \texttt{\tiny{[GM$|$GM]}}
                \texttt{* success: True} \\
\texttt{* lose: False} \\
\texttt{* aborted: False} \\
\texttt{{-}{-}{-}{-}{-}{-}{-}} \\
\texttt{* Shifts: 15.00} \\
\texttt{* Max Shifts: 12.00} \\
\texttt{* Min Shifts: 6.00} \\
\texttt{* End Distance Sum: 0.00} \\
\texttt{* Init Distance Sum: 31.07} \\
\texttt{* Expected Distance Sum: 29.33} \\
\texttt{* Penalties: 12.00} \\
\texttt{* Max Penalties: 16.00} \\
\texttt{* Rounds: 16.00} \\
\texttt{* Max Rounds: 28.00} \\
\texttt{* Object Count: 7.00} \\
            }
        }
    }
    & & \\ \\

\end{supertabular}
}

\end{document}
